\subsubsection{}

\begin{para}
Soient $A$ une $K$-alg\`ebre, $E \in \text{ob}\,D^-(A)$, $\rho : A \longrightarrow A$ un endomorphisme de $K$-alg\`ebres. Notons $A_\rho$ l'anneau $A$ consid\'er\'e comme $A^e$-module \`a gauche par l'homomorphisme $A^e \longrightarrow A$, $a \otimes a' \mapsto \rho(a)a'$. On a un isomorphisme canonique $\Hom_A(E, A_\rho \otimes_A E) \xrightarrow{\sim} \Hom_A(E, E)$ o\`u le second membre est l'ensemble des homomorphismes de $A$-modules \`a gauche $E \longrightarrow E$ qui sont $\rho$-lin\'eaires, i.e.\ tels que $u(ax) = \rho(a)u(x)$. Pour $u \in \Hom_A(E, A_\rho \otimes_A E)$, on notera parfois $\Tr_A^\rho(u)$, ou simplement $\Tr_\rho(u)$, l'\'el\'ement $\Tr_A(u) \in A_K$. Lorsque $\rho$ est un automorphisme int\'erieur, $\rho(a) = gag^{-1}$, on v\'erifie que $\Tr_\rho(u)$ ne d\'epend que de la classe de $g$ dans $A_K^*$ et que, pour $u$ $\rho$-lin\'eaire, on a
\begin{equation}\label{eq:5.13.1}
\Tr_\rho(u) = \text{classe de } g \text{ dans } A_K \cdot \Tr(g^{-1}u)
\end{equation}
o\`u $g^{-1}u$ est l'homomorphisme $A$-lin\'eaire $x \mapsto g^{-1}u(x)$.
\end{para}

\newpage

\subsection{Correspondances \'equivariantes et divisibilit\'e de termes locaux}

\subsubsection{Notations}

\begin{para}
On reprend les notations de I, et on fixe de plus un groupe fini $G$. On se donne des $S$-sch\'emas \`a op\'eration de $G$ (\`a gauche) : $\underline{X}$, $\underline{Y}$, des sous-sch\'emas ferm\'es stables $\underline{A} \subset \underline{X} \times_S \underline{Y}$, $\underline{B} \subset \underline{X} \times_S \underline{Y}$. On note $X$, $Y$, $A$, $B$ les sch\'emas sous-jacents (i.e.\ les quotients par $G$). On suppose que $A$ et $B$ v\'erifient les hypoth\`eses de position g\'en\'erale de I~1.2. On suppose donn\'es de plus des objets $L_X \in \text{ob}\,D(X)$, $L_Y \in \text{ob}\,D(Y)$, munis d'op\'erations de $G$ au-dessus de celles sur $X$ et $Y$ (i.e.\ des homomorphismes $G \longrightarrow \Aut(L_X)$, $G \longrightarrow \Aut(L_Y)$), et des homomorphismes $G$-\'equivariants $u : a_1^*L_X \longrightarrow a_2^!L_Y$, $v : b_2^*L_Y \longrightarrow b_1^!L_X$.
\end{para}

On a alors des termes locaux $\langle u_z, v_z \rangle \in \Lambda$ comme en I~1.1.1. D'autre part, le cup-produit global $\langle \underline{u}, \underline{v} \rangle$ (I~1.1.2) est un \'el\'ement de $H^4_c(X \times_S Y, K_{X \times_S Y}) = \Lambda$, et le th\'eor\`eme 1.2 s'applique pour donner $\langle u_z, v_z \rangle = \langle \underline{u}, \underline{v} \rangle$.

Mais on peut aussi consid\'erer les traces non commutatives des fl\`eches $u_z$ et $v_z$. Plus pr\'ecis\'ement, $(L_X)_x$ et $(L_Y)_y$ sont des complexes parfaits de $\Lambda[G]$-modules, et $u_z$, $v_z$ sont des homomorphismes de $\Lambda[G]$-modules. On peut donc former, d'apr\`es 5.6.3,
\[
\Tr_{\Lambda[G]}(u_z v_z) \in \Lambda[G]_\Lambda \ ,
\]
o\`u $\Lambda[G]_\Lambda$ est le $\Lambda$-module libre sur l'ensemble $G/\sim$ des classes de conjugaison de $G$. De m\^eme, le cup-produit global fournit un \'el\'ement
\[
\langle\!\langle \underline{u}, \underline{v} \rangle\!\rangle \in H^4_c(X \times_S Y, \underline{\Lambda})_{\Lambda[G]}
\]
o\`u $\underline{\Lambda}$ d\'esigne le faisceau constant de valeur $\Lambda[G]_\Lambda$.

\begin{theoreme}[6.2]\label{thm:6.2}
Avec les notations ci-dessus, supposons de plus que $A$ et $B$ soient en position g\'en\'erale (au sens de \emph{I~1.2}). Alors on a
\[
\Tr_{\Lambda[G]}(u_z v_z) = \langle\!\langle \underline{u}, \underline{v} \rangle\!\rangle \ .
\]
En d'autres termes, le terme local non commutatif co\^incide avec le terme global non commutatif.
\end{theoreme}

La d\'emonstration suit le m\^eme sch\'ema que celle de 1.2. On se ram\`ene d'abord, par les m\^emes r\'eductions qu'au $n^{\circ}$~2, au cas o\`u $A$ et $B$ sont des immersions ferm\'ees, et il s'agit de prouver une nullit\'e analogue \`a (2.4.1). Le lemme clef est le suivant :

\begin{lemme}[6.2.1]\label{lem:6.2.1}
Sous les hypoth\`eses de 6.2, supposons que $(L_X)|_U = 0$, $(L_Y)|_V = 0$ (avec les notations de \emph{2.1}). Alors on a
\[
\Tr_{\Lambda[G]}(u_z v_z) = \Tr_{\Lambda[G]}(\underline{u} \cdot \underline{v})
\]
dans $\Lambda[G]_\Lambda$.
\end{lemme}

La d\'emonstration de 6.2.1 est analogue \`a celle de 2.3. On se ram\`ene au cas o\`u $L_X$ et $L_Y$ sont concentr\'es aux points ferm\'es, et on applique la formule de Lefschetz non commutative pour les projections $\{x\} \longrightarrow S$, $\{y\} \longrightarrow S$.

Pour terminer la d\'emonstration de 6.2, on doit encore \'etablir l'analogue non commutatif de 2.5 :

\begin{proposition}[6.2.2]\label{prop:6.2.2}
Sous les hypoth\`eses de 6.2, supposons que $a$ et $b$ sont des immersions ferm\'ees et que $(L_X)_x = 0$, $(L_Y)_y = 0$. D\'esignons par $T$ le localis\'e strict de $X \times_S Y$ en $z = (x,y)$, et identifions $A$, $B$ \`a des sous-sch\'emas ferm\'es de $T$. Alors la fl\`eche de restriction
\[
\Tr_{\Lambda[G]} : R\Gamma_A(T, \RHom(DL_X \otimes L_Y, K_T)) \longrightarrow R\Gamma(B, \RHom(DL_X \otimes L_Y, K_T)|_B)
\]
est nulle dans $D(\Lambda[G]_\Lambda)$.
\end{proposition}

La d\'emonstration occupe le reste du $n^{\circ}$~6. On se ram\`ene, comme au $n^{\circ}$~3, au cas mod\'er\'e, o\`u les repr\'esentations de l'inertie se factorisent \`a travers des $\ell$-groupes. La r\'eduction au cas mod\'er\'e est analogue \`a celle de 3.1--3.3. Il reste alors \`a traiter le cas mod\'er\'e proprement dit, ce qui se fait \`a l'aide d'un th\'eor\`eme de puret\'e (4.3) et d'un calcul explicite des termes locaux.

\begin{para}
Le th\'eor\`eme 6.2 a des cons\'equences importantes pour la divisibilit\'e des termes locaux. Soit $\chi$ une repr\'esentation virtuelle de $G$ sur $\Lambda$ (i.e.\ un \'el\'ement de $R_\Lambda(G)$). On peut consid\'erer le terme local $\chi(u_z, v_z)$ obtenu en appliquant $\chi$ au terme local non commutatif $\Tr_{\Lambda[G]}(u_z v_z)$. Plus pr\'ecis\'ement, si $\chi = \sum n_i [V_i]$, avec $V_i$ des repr\'esentations irr\'eductibles de $G$, on pose
\[
\chi(u_z, v_z) = \sum n_i \Tr(u_z v_z | (L_X)_x \otimes_{\Lambda[G]} V_i) \ .
\]
Alors 6.2 entra\^ine :
\end{para}

\begin{corollaire}[6.3]\label{cor:6.3}
Avec les notations de 6.2, soient $\chi \in R_\Lambda(G)$, et notons $e_\chi = \frac{1}{|G|} \sum_{g \in G} \chi(g^{-1}) g$ l'idempotent central de $\Lambda[G]$ associ\'e \`a $\chi$. Alors
\[
\chi(u_z, v_z) = e_\chi \cdot \langle\!\langle \underline{u}, \underline{v} \rangle\!\rangle \ .
\]
En particulier, si $\chi$ est une repr\'esentation de degr\'e $1$, on a
\[
\chi(u_z, v_z) = \chi(\langle u_z, v_z \rangle^G) \ ,
\]
o\`u $\langle u_z, v_z \rangle^G$ d\'esigne la composante de $\Tr_{\Lambda[G]}(u_z v_z)$ sur la classe de conjugaison triviale.
\end{corollaire}

\begin{para}
Lorsque $G$ est un groupe cyclique, et que l'on prend pour $L_X$, $L_Y$ les faisceaux constants $\Lambda[G]$, et pour $u$, $v$ les correspondances donn\'ees par les multiplications \`a gauche et \`a droite, le th\'eor\`eme 6.2 redonne la formule de divisibilit\'e de Grothendieck (XII~4.5). Plus g\'en\'eralement, pour tout groupe fini $G$ et toute repr\'esentation virtuelle $\chi$, la formule de 6.3 permet de ``diviser canoniquement'' les termes locaux ordinaires attach\'es \`a des correspondances \'equivariantes, comme annonc\'e dans l'introduction.
\end{para}

\vfill

\begin{center}
\rule{0.5\textwidth}{0.4pt}

\textit{Fin de l'Expos\'e III}
\end{center}

\newpage

\section*{Bibliographie compl\'ementaire}

\begin{enumerate}
\item[{[1]}] H.~Cartan, S.~Eilenberg, \textit{Homological Algebra}, Princeton University Press, 1956.

\item[{[2]}] A.~Grothendieck, ``Formule de Lefschetz et rationalit\'e des fonctions $L$'', \textit{S\'eminaire Bourbaki}, expos\'e 279, 1964/65.

\item[{[3]}] J.-L.~Verdier, ``Cat\'egories d\'eriv\'ees (\'etat 0)'', in \textit{SGA $4\frac{1}{2}$}, Springer Lecture Notes in Mathematics 569, 1977.

\item[{[4]}] R.P.~Langlands, ``Modular forms and $\ell$-adic representations'', in \textit{Modular functions of one variable II}, Springer Lecture Notes in Mathematics 349, 1973.

\item[{[5]}] J.~Stallings, ``Centerless groups---an algebraic formulation of Gottlieb's theorem'', \textit{Topology} 4 (1965), 129--134.

\item[{[6]}] M.~Artin, J.-L.~Verdier, ``Refinement of the Thomason--Trobaugh theorem'', (manuscrit).
\end{enumerate}


\clearpage
% ===== Source block: SGA 5 pages 201 250 =====
\phantomsection
\addcontentsline{toc}{section}{SGA 5 pages 201 250}
\section*{SGA 5 pages 201 250}
\setcounter{page}{201}

\selectlanguage{french}


\noindent\textbf{Note:} This document covers pages 201--250 of SGA~5, comprising the end of the preceding Expose (continuation of Section 6 on traces and equivariant cohomology), the Bibliography, and Expose~V (``Systemes projectifs J-adiques'') by J.P.~Jouanolou.

\bigskip
\hrule
\bigskip

\section*{End of preceding Expose (pages 201--215 of PDF)}
\addcontentsline{toc}{section}{End of preceding Expose}

\noindent\textbf{6.7.} Soient $X$ et $c$ comme en 6.5, et $A \longrightarrow B$ un morphisme de $\Lambda$-algebres plates. Vis-a-vis de l'extension et de la restriction des scalaires, les fleches $\Tr_A$ (6.5.3, 6.5.8, 6.5.9) se comportent comme les fleches $\Tr_A$ de 5.6.1, 5.6.2. Le point est que, si $f$ est un $S$-morphisme, les foncteurs $(f^*, f_*, f_!, f^!)$ commutent a l'extension et a la restriction des scalaires, et que par consequent l'etude du comportement de $\Tr_A$ (6.5.3) vis-a-vis de l'extension et de la restriction des scalaires se ramene a celle du comportement de la fleche ``produit tensoriel''
\[
\RHom_A(p_1^{*L}, p_1^{*KA_X}) \overset{L}{\otimes}_{A\,p_2^{*L}} \longrightarrow \RHom_A(p_1^{*L}, p_1^{*KA_X}) \overset{L}{\otimes}_{A\,p_2^{*L}})
\]
et de la fleche ``d'evaluation''
\[
\RHom_A(L, KA_X) \overset{L}{\otimes}_A L \longrightarrow KA_{X,L_\natural}\;.
\]

Ainsi, la commutativite des carres 5.9.1, 5.9.2 entraine celle du carre
\begin{equation}\tag{6.7.1}
\begin{array}{ccc}
\delta^{*c}{}_*\RHom_A(c_1^{*L}, c_2^{!\,L}) & \xrightarrow{\quad\Tr_A\quad} & i_*^c KA_{X^c,L_\natural} \\[8pt]
\Big\downarrow & & \Big\downarrow \\[8pt]
\delta^{*c}{}_*\RHom_B(c_1^{*}(B \overset{L}{\otimes}_A L),\, c_2^{!\,}(B \overset{L}{\otimes}_A L)) & \xrightarrow{\quad\Tr_B\quad} & i_*^c KB_{X^c,L_\natural}
\end{array}\;,
\end{equation}
ou la fleche verticale de gauche est l'extension des scalaires et celle de droite est definie par la fleche canonique $A \overset{L}{\otimes}_{A^e} A \longrightarrow B \overset{L}{\otimes}_{B^e} B$; il en resulte que, pour $u \in \Hom_A(c_1^{*L}, c_2^{!\,L})$, on a
\begin{equation}\tag{6.7.2}
\Tr_B(B \overset{L}{\otimes}_A u) = \varphi\Tr_A(u)\;,
\end{equation}
ou $\varphi : H^0(X^c, KA_{X^c,L_\natural}) \longrightarrow H^0(X^c, KB_{X^c,L_\natural})$ est l'application canonique.

Supposons $A$ de type fini et plat en tant que $\Lambda$-module, $B$ de type fini et plat en tant que $A$-module a gauche, et $A_\natural$, $B_\natural$ plats sur $\Lambda$. On dispose de $\Tr_{B/A} : B_\natural \longrightarrow A_\natural$, on notera encore
\begin{equation}\tag{6.7.3}
\Tr_{B/A} : KB_{X,\natural} \longrightarrow KA_{X,\natural}\;,
\end{equation}
la fleche qui s'en deduit par application de $K_X \overset{L}{\otimes} -$ (6.5.2). Il decoule de 5.10.6 et 5.10.9 que, pour $L \in \mathrm{ob}\, D_{ctf}^-(X)$, on a un carre commutatif
\begin{equation}\tag{6.7.4}
\begin{array}{ccc}
\delta^{*c}{}_*\RHom_B(c_1^{*L}, c_2^{!\,L}) & \xrightarrow{\quad\Tr_B\quad} & i_*^c KB_{X^c,\natural} \\[8pt]
\Big\downarrow & & \Big\downarrow \rlap{$\scriptstyle\Tr_{B/A}$} \\[8pt]
\delta^{*c}{}_*\RHom_A(c_1^{*L}, c_2^{!\,L}) & \xrightarrow{\quad\Tr_A\quad} & i_*^c KA_{X^c,\natural}
\end{array}\;,
\end{equation}
ou la fleche verticale de gauche est la restriction des scalaires. Par suite, pour $u \in \Hom_B(c_1^{*L}, c_2^{!\,L})$, on a, dans $H^0(X^c, KA_{X^c,\natural})$,
\begin{equation}\tag{6.7.5}
\Tr_A(u) = \Tr_{B/A}(\Tr_B(u))\;.
\end{equation}

\medskip
\noindent\textbf{6.8.} Soient $A$, $B$ des $\Lambda$-algebres plates, $X$, $Y$ des $S$-schemas, $Z = X \times_S Y$, $c : C \longrightarrow X^2$, $d : D \longrightarrow Y^2$ des correspondances, $e = c \times_S d : E = C \times_S D \rightarrow Z^2$ le produit. On a donc $Z^e = X^c \times_S Y^d$. Posons $R = A \otimes_\Lambda B$. On definit comme en (III 5.3) le produit tensoriel externe
\begin{equation}\tag{6.8.1}
\Hom_A(c_1^{*L}, c_2^{!\,L}) \otimes \Hom_B(d_1^{*M}, d_2^{!\,M}) \xrightarrow{\overset{L}{\otimes}_S} \Hom_R(e_1^{*}(L \overset{L}{\otimes}_S M), e_2^{!\,}(L \overset{L}{\otimes}_S M))\;,
\end{equation}
et l'on deduit aisement de la commutativite des carres 5.12.1, 5.12.3, que, pour $u \in \Hom_A(c_1^{*L}, c_2^{!\,L})$, $v \in \Hom_B(d_1^{*M}, d_2^{!\,M})$, on a, dans $H^0(Z^e, KR_{Z^c,L_\natural})$,
\begin{equation}\tag{6.8.2}
\Tr_R(u \overset{L}{\otimes}_S v) = \Tr_A(u)\,\Tr_B(v)\;,
\end{equation}
le produit au second membre etant defini par la fleche canonique
\begin{equation}\tag{6.8.3}
H^0(X^c, KA_{X^c,L_\natural}) \otimes H^0(Y^d, KB_{Y^d,L_\natural}) \longrightarrow H^0(Z^c; KR_{Z^c,L_\natural})\;.
\end{equation}

\medskip
\noindent\textbf{6.9.} Soit $G$ un groupe fini. Si $X$ est un $G$-$S$-schema (i.e.\ un $S$-schema muni d'une action de $G$ par $S$-automorphismes$^{(*)}$), et $A$ une $\Lambda$-algebre$^{(**)}$ nous noterons $D(G,A^X)$ la categorie derivee de celle des faisceaux de $G$-$A$-modules a gauche sur $X$ (pour les notions de $G$-faisceau d'ensembles, d'anneaux, de modules, etc.) sur $X$ (pour la topologie etale), nous renvoyons le lecteur a (XV 1)), et nous designerons par l'indice $c$ (resp.\ $ctf$) la sous-categorie pleine formee des complexes de $G$-$A$-modules qui, en tant que complexes de $A$-modules, sont a cohomologie constructible (resp.\ de tor-dimension finie et a cohomologie constructible).

De meme, si $B$ est une $\Lambda$-algebre, nous noterons $D(G,A_B^X)$ la categorie derivee de celle des $G$-$(A,B)$-bimodules sur $X$. Si $G$ agit trivialement sur $X$, un $G$-$A$-module a gauche sur $X$ n'est autre qu'un $A[G]$-module a gauche sur $X$ (sans groupe d'operateurs), donc $D(G,A^X) = D(A[G]^X)$; comme $G$ est fini, on a de plus $D_c(A^X) = D_c(A[G]^X)$, mais on a seulement une inclusion $D_{ctf}(A[G]^X) \subset D_{ctf}(G,A^X)$.

Le formalisme de dualite de (SGA 4 XVII, XVIII) s'etend sans peine au cas des $G$-$S$-schemas. Comme le foncteur d'oubli de l'action de $G$ est fidele, il suffit d'etendre de facon naturelle la definition des ``six operations'' ($\overset{L}{\otimes}$, $\RHom$, $f^*$, $f_*$, $f_!$, $f^!$) et des fleches canoniques envisagees en (loc.\ cit) (changement de base, Kunneth, dualite, etc.). Cette extension n'offre pas de difficulte (pour la definition des foncteurs $f_!$, $f^!$, il suffit d'utiliser des compactifications equivariantes). Nous laissons au lecteur le soin de paraphraser les definitions et resultats de 6.2 a 6.8 dans le cadre des $G$-$S$-schemas. Si, pour $i=1,2$, $X_i$ est un $G$-$S$-schema, $c : C \longrightarrow X = X_1 \times_S X_2$ un $G$-$S$-morphisme, et $L_i \in \mathrm{ob}\, D_{ctf}(G, A_{X_i}^{X_i})$, les elements de $\Hom_A(c_1^{*L_1}, c_2^{!\,L_2})$ ($\Hom_A = \Hom$ dans la categorie $D(G, A^-)$) s'appelleront \underline{correspondances (cohomologiques) equivariantes de $L_1$ a $L_2$ a support dans $c$}. On notera que, lorsque $G$ agit trivialement sur $X_i$ ($i=1,2$) et que $L_i \in \mathrm{ob}\, D_{ctf}(A[G]^{X_i})$, l'isomorphisme 6.4.1 (dans $D(G,X)$) se ``raffine'' en un isomorphisme
\[
D_{A[G]}L_1 \overset{L}{\otimes}_{S,A[G]} L_2 \xrightarrow{\;\sim\;} \RHom_{A[G]}(p_1^{*L_1}, p_2^{!\,L_2})\;.
\]

Il en resulte que, dans une situation $G$-equivariante de type 6.5, avec action triviale de $G$ sur les espaces, et $L \in \mathrm{ob}\, D_{ctf}(A[G]^X) \subset D_{ctf}(G, A^X)$) on peut, pour $u \in \Hom_{A[G]}(c_1^{*L}, c_2^{!\,L})$, definir un terme local $\Tr_{A[G]}(u) \in H^0(X^c, KA[G]_{X^c,L_\natural})$, plus fin (en un sens evident) que le terme local $\Tr_A(u) \in H^0(G, X^c, KA_{X^c,L_\natural})$ obtenu en utilisant seulement la tor-dimension finie de $L$ par rapport a $A$.

\medskip
\noindent\textbf{6.10.} Soient $X$, $c$ comme en 6.5, et $L \in \mathrm{ob}\, D_{ctf}(A[G]^X)$. Pour $u \in \Hom_{A[G]}(c_1^{*L}, c_2^{!\,L})$, on dispose, par 6.5.9, d'une trace locale
\begin{equation}\tag{6.10.1}
\Tr_{A[G]}(u) \in H^0(X^c, KA[G]_{X^c,\natural}) = H^0(X^c; K_{X^c} \otimes A[G_\natural])\;,
\end{equation}
ou $G_\natural$ est l'ensemble des classes de conjugaison de $G$, et pour tout $x \in G_\natural$, on peut considerer la ``valeur de $\Tr_{A[G]}(u)$ en $x$'', $\Tr_{A[G]}(u)(x) \in H^0(X^c, K_{X^c})$, image de $\Tr_{A[G]}(u)$ par la projection $H^0(X^c, K_{X^c} \otimes A[G_\natural]) \longrightarrow H^0(X^c, K_{X^c})$ definie par $x$. Il resulte de 6.7.5 et 5.11.2 que l'on a
\begin{equation}\tag{6.10.2}
\Tr_A(u) = |G|\,\Tr_{A[G]}(u)(e)\;,
\end{equation}
ou $e$ designe l'element neutre de $G$. Plus generalement, soit $g \in G$, notons $Z_g$ le centralisateur de $g$, et $gu \in \Hom_{A[Z_g]}(c_1^{*L}, c_2^{!\,L})$ la correspondance composee de $u$ avec la multiplication a gauche par $g : c_1^{*L} \longrightarrow c_1^{*L}$. On a $L \in \mathrm{ob}\, D_{ctf}(A[Z_g]^X)$, et par suite 6.10.2 entraine
\begin{equation}\tag{6.10.3}
\Tr_A(gu) = |Z_g|\,\Tr_{A[Z_g]}(gu)(e)\;.
\end{equation}
Par le meme raisonnement qu'en 5.11.3, on deduit de 6.7.4 la formule
\begin{equation}\tag{6.10.4}
\Tr_{A[Z_g]}(gu)(e) = \Tr_{A[G]}(u)(g^{-1})\;,
\end{equation}
qui permet, si on le desire, de reecrire 6.10.3 sous la forme
\begin{equation}\tag{6.10.5}
\Tr_A(gu) = |Z_g|\,\Tr_{A[G]}(u)(g^{-1})\;.
\end{equation}

\medskip
\noindent\textbf{6.11.} Soient $f : X \longrightarrow Y$ un $G$-$S$-morphisme propre, et
\begin{equation}\tag{6.11.1}
\begin{array}{ccc}
X \times_S X & \xleftarrow{\quad c\quad} & C \\[4pt]
\llap{$f \times_S f$}\Big\downarrow\rlap{$\quad$} & & \Big\downarrow \\[8pt]
Y \times_S Y & \xleftarrow{\quad d\quad} & D
\end{array}
\end{equation}
un carre commutatif $G$-equivariant de $G$-$S$-schemas, avec $c$ et $d$ propres. On suppose que $G$ agit trivialement sur $Y$ et $D$. On note
\begin{equation}\tag{6.11.2}
f^{c/d}\;\text{(ou }f^c) : X^c \longrightarrow Y^d
\end{equation}
le morphisme induit par 6.11.1 sur les schemas des points fixes de $c$ et $d$, avec les notations de 6.5.1. Soient $L \in \mathrm{ob}\, D_{ctf}(G,X)$, et $u \in \Hom(c_1^{*L}, c_2^{!\,L})$ une correspondance $G$-equivariante a support dans $c$, d'ou, par image directe, une correspondance $G$-equivariante $f_*u \in \Hom_{A[G]}(d_1^{*f_*L}, d_2^{!\,f_*L})$. Pour $g \in G$, nous noterons $Z_g$ le centralisateur de $g$, et
\begin{equation}\tag{6.11.3}
gu \in \Hom((gc)_1^{*L}, c_2^{!\,L})
\end{equation}
la correspondance $Z_g$-equivariante a support dans
\begin{equation}\tag{6.11.4}
gc \overset{\mathrm{dfn}}{=} (gc_1 = c_1g, c_2) : C \longrightarrow X \times_S X
\end{equation}
definie comme composee des fleches
\[
g^*(c_1^{*L}) \xrightarrow{\quad\alpha_g\quad} c_1^{*L} \xrightarrow{\quad u\quad} c_2^{!\,L}\;,
\]
ou $\alpha_g$ est la fleche donnant l'action de $G$ sur $c_1^{*L}$. Il est clair que l'on a
\begin{equation}\tag{6.11.5}
f_*(gu) = g(f_* u)\;.
\end{equation}

\medskip
\noindent\textbf{Proposition 6.12} --- \underline{Sous les hypotheses de 6.11, on suppose que l'on a} $f_* L \in \mathrm{ob}\, D_{ctf}(A[G]^Y)$. \underline{Alors, pour tout} $g \in G$, \underline{on a}
\begin{equation}\tag{6.12.1}
f_*^{gc}\Tr_A(gu) = |Z_g|\,\Tr_{A[G]}(f_* u)(g^{-1})\;,
\end{equation}
\underline{ou} $f_*^{gc} : H^0(X^{gc}, K_{X^{gc}}) \longrightarrow H^0(Y^d, K_{Y^d})$ \underline{est le morphisme ``de Gysin'' defini par la fleche d'adjonction} $f_*^{gc}K_{X^{gc}} = f_*^{gc}f^{gc!}K_{Y^d} \longrightarrow K_{Y^d}$.

D'apres 6.10.5 applique a $f_* u$, et 6.11.5, on a en effet
\[
\Tr_A(f_*(gu)) = |Z_g|\,\Tr_{A[G]}(f_* u)(g^{-1})\;,
\]
et par suite 6.12.1 resulte de la formule de Lefschetz (III 4.5.1).

\medskip
\noindent\textbf{6.13.} Soit $\ell$ un nombre premier $\not= p$. Pour tout entier $n \geq 1$, posons $\Lambda_n = \bZ/\ell^n$. Sous les hypotheses de 6.11, soit $L = (L_n)$ un systeme projectif d'objets de $D_{ctf}(G,X,\Lambda_n)$ tel que $L_m \overset{L}{\otimes}_{\Lambda_m} \Lambda_n \xrightarrow{\sim} L_n$ pour $m \geq n$, et soit $u \in \Hom(c_1^{*L}, c_2^{!\,L})$ un systeme projectif de correspondances equivariantes $u_n \in \Hom(c_1^{*L_n}, c_2^{!\,L_n})$ tel que, pour $m \geq n$, $u_m \overset{L}{\otimes}_{\Lambda_m} \Lambda_n = u_n$ modulo l'isomorphisme $L_m \overset{L}{\otimes}_{\Lambda_m} \Lambda_n = L_n$. En vertu de la compatibilite des traces locales a l'extension des scalaires, les termes locaux (pour $g \in G$) $\Tr_{\Lambda_n}(gu_n) \in H^0(X^{gc}, K\Lambda_{n,X^{gc}})$ definissent un element
\[
\Tr_{\bZ_\ell}(gu) \in H^0(X^{gc}, K_{\ell,X^{gc}}) \overset{\mathrm{dfn}}{=} \varprojlim_n H^0(X^{gc}, K\Lambda_{n,X^{gc}})\;,
\]
et les $f_*^{gc}\Tr_{\Lambda_n}(gu_n) = \Tr_{\Lambda_n}(gf_* u_n)$ un element
\[
f_*\Tr_{\bZ_\ell}(gu) \in H^0(Y^d, K_{\ell,Y^d}) \overset{\mathrm{dfn}}{=} \varprojlim_n H^0(Y^d, K\Lambda_{n,Y^d})\;.
\]
Supposons de plus que $f_*^{M_n} \in \mathrm{ob}\, D_{ctf}(\Lambda_n[G]^Y)$ pour tout $n$. Alors les traces locales $\Tr_{\Lambda_n[G]}(f_* u_n)$ definissent de meme un element
\[
\Tr_{\bZ_\ell[G]}(f_* u) \in H^0(Y^d, K\bZ_\ell[G]_{Y^d,\natural}) \overset{\mathrm{dfn}}{=} \varprojlim_n H^0(Y^d, K\Lambda_n[G]_{Y^d,\natural})\;.
\]
Il decoule de 6.12 que l'on a
\begin{equation}\tag{6.13.1}
f_*^{gc}\Tr_{\bZ_\ell}(gu) = |Z_g|\,\Tr_{\bZ_\ell[G]}(f_* u)(g^{-1})\;.
\end{equation}

\medskip
\noindent\textbf{6.14.} Soient $X$ un $G$-$S$-schema, $j : U \longrightarrow X$ une immersion ouverte equivariante, $c : C \longrightarrow X \times_S X$ un morphisme de $G$-$S$-schemas tel que $c_2$ ($= p_2c$) : $C \longrightarrow X$ soit quasi-fini et de tor-dimension finie. On suppose que $c_1^{-1}(U) \subset c_2^{-1}(U)$; on a donc des carres commutatifs (6.14.1).

Le morphisme trace $\Tr_{c_2} : c_{2!}\Lambda_C \longrightarrow \Lambda_X$ definit alors une correspondance $G$-equivariante filtree
\begin{equation}\tag{6.14.2}
\underline{c} \in \Hom(c_1^*\Lambda_X, c_2^{!\,}\Lambda_X)\;,
\end{equation}
$\Lambda_X$ etant muni de la filtration definie par le sous-module $i_!\Lambda_U$ (i.e.\ $F^n\Lambda_X = \Lambda_X$ pour $n \leq 0$, $F^1\Lambda_X = i_!\Lambda_U$, $F^n\Lambda_X = 0$ pour $n \geq 1$). Le gradue associe se compose des deux correspondances $G$-equivariantes
\begin{align*}
\mathrm{gr}^0\underline{c} & = \underline{c}_{X-U,X} \in \Hom(c_1^*(\Lambda_{X-U,X}), c_2^{!\,}(\Lambda_{X-U,X}))\;, \tag{6.14.3} \\
\mathrm{gr}^1\underline{c} & = \underline{c}_{U,X} \in \Hom(c_1^*(\Lambda_{U,X}), c_2^{!\,}(\Lambda_{X-U,X}))\;,
\end{align*}
ou l'on a pose $\Lambda_{X-U,X} = j_*\Lambda_{X-U}$, $\Lambda_{U,X} = i_!\Lambda_U$. On a
\begin{equation}\tag{6.14.4}
\underline{c}_{X-U,U} = j_*\underline{c}_{X-U}\;,
\end{equation}
ou $\underline{c}_{X-U} \in \Hom(c_1^*\Lambda_{X-U}, c_2^{!\,}\Lambda_{X-U})$ est la correspondance definie par $\Tr_{c_2,X-U}$. Nous ecrirons $\underline{c}$, $\underline{c}_{U,X}$, etc., quand nous voudrons preciser l'anneau de base.

Oublions l'action de $G$. On a des termes locaux
\begin{equation}\tag{6.14.5}
\Tr_\Lambda(\underline{c}),\; \Tr_\Lambda(\underline{c}_{U,X}),\; \Tr_\Lambda(\underline{c}_{X-U,X}) \in H^0(X^c, K_{X^c})\;,
\end{equation}
qui, en vertu de l'additivite des traces (III 4.13.1), verifient la relation
\begin{equation}\tag{6.14.6}
\Tr_\Lambda(\underline{c}) = \Tr_\Lambda(\underline{c}_{U,X}) + \Tr_\Lambda(\underline{c}_{X-U,X})\;.
\end{equation}
De plus, d'apres 6.14.4, on a
\begin{equation}\tag{6.14.7}
\Tr_\Lambda(\underline{c}_{X-U,X}) = j_*^c\Tr_\Lambda(\underline{c}_{X-U})\;,
\end{equation}
ou $j_*^c : H^0((X-U)^c, K_{(X-U)^c}) \longrightarrow H^0(X^c, K_{X^c})$ est le morphisme image directe definie par $j^c : (X-U)^c \longrightarrow X^c$. Rappelons (III 4.3) que si $X$ est lisse, $c$ une immersion fermee et $X^c$ fini, $\Tr_\Lambda(\underline{c})$ est l'image dans $\Lambda$ de la multiplicite d'intersection $(C.X)$ de $C$ avec la diagonale; si de plus $X-U$ est lisse, $\Tr_\Lambda(\underline{c}_{X-U})$ est l'image dans $\Lambda$ de $(c_2^{-1}(X-U).(X-U))$. Notons aussi que, si $X-U$ est propre sur $S$, et $c|X-U$ est la diagonale, alors la formule de Lefschetz (III 4.7) entraine que $\int_{X-U} \Tr(\underline{c}_{X-U}) \in \Lambda$ est la caracteristique d'Euler-Poincare de $X-U$ (a valeurs dans $\Lambda$).

\medskip
\noindent\textbf{6.15.} Soient un diagramme $G$-equivariant 6.11.1 (avec $f$, $c$, $d$ propres), et $U \longrightarrow X$ une immersion ouverte $G$-equivariante. On suppose que $c_2$ est quasi-fini et de tor-dimension finie, $c_1^{-1}(U) \subset c_2^{-1}(U)$, que $G$ agit trivialement sur $C$ et $D$, et que $f|U : U \longrightarrow f(U) = V$ est un revetement etale galoisien de groupe $G$. Cette derniere hypothese implique que $Rf_*(\Lambda_{U,X}) \in \mathrm{ob}\, D_{ctf}(A[G]^Y)$ : en fait $Rf_*(\Lambda_{U,X})$ est concentre en degre $0$ et prolonge par zero d'un faisceau sur $V$ localement isomorphe a $\Lambda[G]$. On peut donc appliquer 6.12 a $L = \Lambda_{U,X}$, $u = \underline{c}_{U,X}$, et, compte tenu de 6.14.6, 6.14.7, on obtient la

\medskip
\noindent\textbf{Proposition 6.16} --- \underline{Sous les hypotheses de 6.15, on a, pour tout} $g \in G$,
\begin{equation}\tag{6.16.1}
f_*^{gc}\Tr_\Lambda(\underline{gc}) - (fj)_*^{gc}\Tr_\Lambda(\underline{gc}_{X-U}) = |Z_g|\,\Tr_{A[G]}(f_*\underline{c}_{U,X})(g^{-1})\;,
\end{equation}
\underline{ou} $j : X-U \longrightarrow X$ \underline{est l'inclusion, et} $Z_g$ \underline{designe le centralisateur de} $g$.

\medskip
Appliquant 6.16 au systeme projectif de correspondances $\underline{c}_\ell = (\underline{c}^{\Lambda_n})$, ou $\Lambda_n = \bZ/\ell^n$ ($\ell$ premier $\not= p$), et tenant compte de ce que 6.14.6, 6.14.7 donnent, par passage a la limite et application de $f_*^{gc}$,
\[
f_*^{gc}\Tr_{\bZ_\ell}(\underline{gc}_\ell) = f_*^{gc}\Tr_{\bZ_\ell}(\underline{gc}_{\ell,U,X}) + (fj)_*^{gc}\Tr_{\bZ_\ell}(\underline{gc}_{\ell,X-U})\;,
\]
on en deduit, d'apres 6.13,

\medskip
\noindent\textbf{Corollaire 6.17} --- \underline{Sous les hypotheses de 6.15, on a, pour tout} $g \in G$,
\begin{equation}\tag{6.17.1}
f_*^{gc}\Tr_{\bZ_\ell}(\underline{gc}_\ell) - (fj)_*^{gc}\Tr_{\bZ_\ell}(\underline{gc}_{\ell,X-U}) = |Z_g|\,\Tr_{\bZ_\ell[G]}(f_*\underline{c}_{\ell,U,X})(g^{-1})\;,
\end{equation}
\underline{avec les notations de} 6.16.

\medskip
Compte tenu de la compatibilite des traces locales a la localisation et des remarques faites en 6.14, on voit que 6.17 resout affirmativement, dans une formulation plus generale, la conjecture de divisibilite de Grothendieck (XII 4.5), dans le cas ou l'endomorphisme $\varphi$ de $G$ considere dans (loc.\ cit.) est l'identite.

En fait, il n'est pas difficile, en s'inspirant de 5.13, de donner des enonces de divisibilite analogues a 6.12, 6.13, 6.16, 6.17, dans le cas ou $c$ n'est pas $G$-equivariante, mais telle que $c_2$ le soit et que $c_1$ verifie $c_1(xg) = c_1(x)\varphi(g)$, pour tout point $x$ de $G$, $\varphi$ etant un endomorphisme donne de $G$: les formules de divisibilite s'ecrivent comme plus haut, avec $Z_g$ remplace par le $\varphi$-centralisateur de $g^{-1}$ (5.13.5).

\medskip
\noindent\textbf{6.18.} Posons
\begin{equation}\tag{6.18.1}
N_G = \sum_{g \in G} g \in \Lambda[G]\;.
\end{equation}
Soit $X$ un $S$-schema. Pour $L \in \mathrm{ob}\, D(A[G]^X)$, la multiplication a gauche par $N_G$ induit une fleche de $D(X)$
\begin{equation}\tag{6.18.2}
N_G : \Lambda \overset{L}{\otimes}_{A[G]} L \longrightarrow \RHom_{A[G]}(\Lambda, L)\;,
\end{equation}
ou $\Lambda$ est considere comme $\Lambda[G]$-algebre par l'augmentation naturelle (envoyant tout $g \in G$ sur $1$).

\medskip
\noindent\textbf{Lemme 6.18.3} --- \underline{Si} $L \in \mathrm{ob}\, D_{ctf}(A[G]^X)$, \underline{la fleche} 6.18.2 \underline{est un isomorphisme.}

Par devissage on se ramene en effet a supposer $L$ parfait, puis $L = \Lambda[G]$, et l'assertion decoule de ce que $N_G$ (6.18.1) est une base du module des invariants sous $G$ du $\Lambda[G]$-module a gauche $\Lambda[G]$.

\medskip
\noindent\textbf{Proposition 6.19} --- \underline{Dans la situation de} 6.11, \underline{soient} $M \in \mathrm{ob}\, D_{ctf}(Y)$, $v \in \Hom(d_1^{*M}, d_2^{!\,M})$, $a : M \longrightarrow f_* L$ \underline{une fleche de} $D(A[G]^Y)$ ($M$ etant regarde comme objet de $D(A[G]^Y)$ par l'augmentation $\Lambda[G] \longrightarrow \Lambda$) \underline{tels que le carre}
\[
\begin{array}{ccc}
d_1^{*M} & \xrightarrow{\quad v\quad} & d_2^{!\,M} \\[6pt]
\llap{$d_1^{*a}$}\Big\downarrow\rlap{$\quad$} & & \Big\downarrow\rlap{$\scriptstyle d_2^{!\,a}$} \\[8pt]
d_1^{*f_* L} & \xrightarrow{\quad f_* u\quad} & d_2^{!\,f_* L}
\end{array}
\]
\underline{soit commutatif. On suppose que} $f_* M \in \mathrm{ob}\, D_{ctf}(A[G]^Y)$ \underline{et que} $a$ \underline{induit un isomorphisme} $M \xrightarrow{\sim} \RHom_{A[G]}(\Lambda, f_* L)$. \underline{Alors on a}
\begin{equation}\tag{6.19.1}
\Tr_\Lambda(v) = \sum_{g \in G_\natural} \Tr_{A[G]}(f_* u)(g^{-1})\;.
\end{equation}

D'apres 6.18.3, $a$, compose avec $N_G^{-1}$, donne un isomorphisme $M \xrightarrow{\sim} \Lambda \overset{L}{\otimes}_{A[G]} f_* L$, par lequel $v$ s'identifie a $\Lambda \overset{L}{\otimes}_{A[G]} f_* u$. D'apres 6.7.2, $\Tr_\Lambda(v)$ est donc l'image de $\Tr_{A[G]}(f_* u)$ par la fleche induite par l'augmentation $\Lambda[G] \longrightarrow \Lambda$; or celle-ci, par definition, envoie chaque classe de conjugaison sur $1$, d'ou la conclusion.


\medskip
\noindent\textbf{6.20.} Soient, pour $i = 1,2$, $X_i$ un $G$-$S$-schema, $X = X_1 \times_S X_2$, $c : C \longrightarrow X$ un $G$-$S$-morphisme tel que $c_2$ soit quasi-fini et de tor-dimension finie. Soient $L_i \in \mathrm{ob}\, D^+(G,X_i)$. Le morphisme trace $\Tr_{c_2} : c_{2!}c_2^{*L_2} \longrightarrow L_2$ definit par adjonction une fleche (de $D(G,X)$) notee encore
\begin{equation}\tag{6.20.1}
\Tr_{c_2} : c_2^{*L_2} \longrightarrow c_2^{!\,L_2}\;,
\end{equation}
qui permet d'associer a tout morphisme ($G$-equivariant) $u \in \Hom(c_1^{*L_1}, c_2^{*L_2})$ une correspondance ($G$-equivariante)
\begin{equation}\tag{6.20.2}
u^! = \Tr_{c_2} \circ u \in \Hom(c_1^{*L_1}, c_2^{!\,L_2})\;.
\end{equation}

\medskip
\noindent\textbf{Proposition 6.21} --- \underline{Sous les hypotheses de 6.11, on suppose donne de plus une immersion ouverte} $i : V \hookrightarrow Y$ \underline{telle que} $(fc_1)^{-1}(V)$ \underline{soit connexe, que} $d_1^{-1}(V) \subset d_2^{-1}(V)$, \underline{et que} $f$ \underline{induise un revetement principal galoisien de groupe} $G$, $U = f^{-1}(V) \xrightarrow{f_V} V$. \underline{On suppose en outre que} $M$ \underline{est de la forme} $M = i_!E$, \underline{ou} $E \in \mathrm{ob}\, D(V)$, \underline{et qu'on a un isomorphisme} $\alpha : f_V^*E \xrightarrow{\sim} \Lambda_X \otimes_\Lambda E_0$ \underline{de} $D(G,X)$, \underline{avec} $E_0 \in \mathrm{ob}\, D(\Lambda[G])$, \underline{parfait sur} $\Lambda$. \underline{On se donne} $v \in \Hom(d_1^{*M}, d_2^{*M})$ \underline{comme en} 6.21, \underline{et l'on note} $v_0 \in \Hom_{\Lambda[G]}(E_0, E_0)$ \underline{l'homomorphisme defini, grace a} $\alpha$, \underline{par} $f^* v|c_1^{-1}(U)$. \underline{Alors, si} $\underline{c}_{U,X}$ \underline{designe la correspondance definie en} 6.14.3, \underline{on a} (\underline{avec la notation} 6.20.2)
\begin{equation}\tag{6.23.1}
\Tr_\Lambda(v^!) = \sum_{g \in G_\natural} \Tr_{A[G]}(f_*\underline{c}_{U,X})(g^{-1})\,\Tr_\Lambda(gv_0)\;.
\end{equation}

(Noter que, grace au caractere central de $\Tr_\Lambda$, le second membre ne depend pas du choix de $\alpha$.)

Par la formule de projection $f_* f^*\Lambda_Y \otimes M \xrightarrow{\sim} f_* f^* M$ et 6.15, $M$ verifie les hypotheses de 6.22. D'autre part, $\alpha$ definit un isomorphisme $f^* M \xrightarrow{\sim} \Lambda_{U,X} \otimes E_0$ par lequel $(f^* v)^!$ s'identifie au produit tensoriel externe $\underline{c}_{U,X} \otimes v_0$. Par la formule de projection, $f_* f^* M$ s'identifie donc a $f_*(\Lambda_{U,X}) \otimes E_0$, et $f_*(f^* v)^!$ au produit tensoriel externe $f_*\underline{c}_{U,X} \otimes v_0$. Pour $g \in G$, la correspondance ($Z_g$-equivariante) $gf_*(f^* v)^!$ s'identifie donc a $gf_*\underline{c}_{U,X} \otimes gv_0$. Comme, d'apres 6.10.4, on a
\[
\Tr_{A[G]}(f_*(f^* v)^!)(g^{-1}) = \Tr_{A[Z_g]}(gf_*(f^* v)^!)(e)\;,
\]
il suffit, compte tenu de 6.22.1, de prouver la formule
\[
\Tr_{A[Z_g]}(gf_*\underline{c}_{U,X} \otimes gv_0)(e) = \Tr_{A[G]}(f_*\underline{c}_{U,X})(g^{-1})\,\Tr_\Lambda(gv_0)\;,
\]
ou, ce qui revient au meme, d'apres 6.10.4,
\begin{equation}\tag{6.23.2}
\Tr_{A[Z_g]}(gf_*\underline{c}_{U,X} \otimes gv_0)(e) = \Tr_{A[Z_g]}(gf_*\underline{c}_{U,X})(e)\,\Tr_\Lambda(gv_0)\;.
\end{equation}
Quitte a remplacer $G$ par $Z_g$, on peut supposer que $g=e$, et, revenant a la definition des traces, on voit qu'il suffit d'etablir le lemme suivant :

\medskip
\noindent\textbf{Lemme 6.23.3} --- (cf.\ (SGA 4$^{1/2}$ Rapport 4.7)). \underline{Soient} $P$ \underline{un} $G_V$-\underline{module localement facteur direct d'un module libre de type fini,} $P' = i_!P$, \underline{et} $Q$ \underline{un complexe de} $\Lambda[G]$-\underline{modules, parfait sur} $\Lambda$. \underline{Alors} $P \overset{L}{\otimes}_\Lambda Q$, \underline{considere comme objet de} $D(\Lambda[G]^V)$ \underline{par l'action diagonale de} $G$, \underline{est parfait sur} $\Lambda[G]$, \underline{et l'on a un carre commutatif}
\[
\begin{array}{ccc}
\delta^*\RHom_{\Lambda[G]}(p_1^{*P'}, p_2^{!\,P'}) \overset{L}{\otimes} \RHom_{\Lambda[G]}(Q,Q) & \xrightarrow{\quad(1)\quad} & K_Y \\[8pt]
\llap{$(*)$}\Big\downarrow\rlap{\quad(2)} & & \parallel \\[10pt]
\delta^*\RHom_{\Lambda[G]}(p_1^{*P'} \overset{L}{\otimes} Q,\, p_2^{!\,P'} \overset{L}{\otimes} Q) & \xrightarrow{\quad(3)\quad} & K_Y
\end{array}\;,
\]
ou $\delta : Y \longrightarrow Y \times_S Y$ \underline{est la diagonale,} (1) \underline{est produit tensoriel des fleches}
\[
\delta^*\RHom_{\Lambda[G]}(p_1^{*P'}, p_2^{!\,P'}) \xrightarrow{\;6.5.3\;} K\Lambda[G]_{Y,\natural} \xrightarrow{\;x \mapsto x(e)\;} K_Y
\]
\underline{et}
\[
\RHom_{\Lambda[G]}(Q,Q) \xrightarrow{\;\Tr_\Lambda\;} \Lambda\;,
\]
(2) \underline{est la fleche deduite du produit tensoriel externe par restriction selon la diagonale} $\Lambda[G] \longrightarrow \Lambda[G] \otimes \Lambda[G]$, \underline{et} (3) \underline{est compose de} 6.5.3 (\underline{definie car} $P' \overset{L}{\otimes} Q \in \mathrm{ob}\, D_{ctf}(\Lambda[G]^Y)$) \underline{et de} $x \mapsto x(e)$.

La premiere assertion est evidente, car on se ramene a $P = \Lambda[G]_V$, et $P \overset{L}{\otimes}_\Lambda Q$ est alors isomorphe au produit tensoriel de $P$ par $Q$ considere comme muni de l'action triviale de $G$. Pour la seconde, notons d'abord que, d'apres 6.5.4, 6.5.5, on a
\[
\delta^*\RHom_{\Lambda[G]}(p_1^{*P'}, p_2^{!\,P'}) \xrightarrow{\sim} D_{\Lambda[G]}(i_!P) \overset{L}{\otimes}_{\Lambda[G]} i_!P \xrightarrow{\sim} i_!(D_{\Lambda[G]}P \overset{L}{\otimes}_{\Lambda[G]} P)\;,
\]
de sorte qu'on est ramene a verifier que le carre induit par $(*)$ sur $V$ commute. On peut donc supposer que $V=Y$. Le carre $(*)$ s'ecrit alors
\[
\begin{array}{ccc}
\RHom_{\Lambda[G]}(P, K_Y \overset{L}{\otimes} P) \overset{L}{\otimes} \RHom_{\Lambda[G]}(Q,Q) & \xrightarrow{\quad(1)\quad} & K_Y \\[10pt]
\llap{$(**)$}\Big\downarrow & & \parallel \\[12pt]
\RHom_{\Lambda[G]}(P \overset{L}{\otimes} Q, K_Y \overset{L}{\otimes} P \overset{L}{\otimes} Q) & \xrightarrow{\quad(3)\quad} & K_Y
\end{array}\;,
\]
ou (2) est le produit tensoriel externe, (1) est donne par $\Tr_{\Lambda[G]}(-)(e) \otimes \Tr_\Lambda$ (traces au sens du n\textdegree\,5), (3) par $\Tr_{\Lambda[G]}(-)(e)$. Par descente, on se ramene a $P = \Lambda[G]_Y$, et, par la remarque du debut, on peut alors supposer que $G$ agit trivialement sur $Q$. L'assertion est alors consequence immediate de 5.12.5.

\medskip
\noindent\textbf{6.24.} Si $\Lambda$ est annule par une puissance d'un nombre premier $\ell \not= p$, on peut reecrire le second membre de 6.23.1 sous la forme
\[
\sum_{g \in G_\natural} \Tr_{\bZ_\ell[G]}(f_*\underline{c}_{\ell,U,X})(g^{-1})\,\Tr_\Lambda(gv_0)\;,
\]
avec les notations de 6.17. La formule 6.17.1 montre que les termes locaux definis par Grothendieck dans (XII 6.2) coincident avec les termes locaux de Verdier, au moins lorsque l'endomorphisme $\varphi$ de $G$ considere dans (loc.\ cit.) est l'identite (mais, comme deja observe, cette restriction n'est pas serieuse). Si $Y$ est propre sur $S$, on obtient, en conjuguant 6.23.1 avec la formule de Lefschetz pour $Y \longrightarrow S$ (III 4.7), une formule generalisant (XII 6.3) (et a fortiori la formule de Lefschetz pour Frobenius demontree dans (SGA 4$^{1/2}$ Rapport) et la Partie I du present expose).

\bigskip
\hrule
\bigskip

\section*{BIBLIOGRAPHIE}
\addcontentsline{toc}{section}{Bibliographie}

\begin{enumerate}
\item[{[1]}] Cartan, H.\ et Eilenberg, S.\ --- Homological Algebra, Princeton University Press, 1956.

\item[{[2]}] Grothendieck, A.\ --- Formule de Lefschetz et rationalite des fonctions $L$, Seminaire Bourbaki 64/65, n\textdegree\,279, in Dix exposes sur la cohomologie des schemas, North Holland, Pub.\ Co., 1968.

\item[{[3]}] Illusie, L.\ --- Complexe cotangent et deformations I, II, Lecture Notes in Mathematics 239, 283, Springer-Verlag, 1971, 1972.

\item[{[4]}] Langlands, R.P.\ --- Modular forms and $\ell$-adic representations, in Modular Functions of One Variable II, Lecture Notes in Mathematics 349, Springer-Verlag, 1973.

\item[{[5]}] Stallings, J.\ --- Centerless groups -- an algebraic formulation of Gottlieb's theorem, Topology, 4, p.\ 129--134, 1965.

\item[{[6]}] Verdier, J.-L.\ --- The Lefschetz fixed point formula in etale cohomology, in Proceedings of a conference on local fields, Springer-Verlag, 1967.
\end{enumerate}

\bigskip
\noindent SGA 4$^{1/2}$ : Cohomologie etale, par P.\ Deligne, Lecture Notes in Mathematics 569, Springer-Verlag 1977.

\medskip
\noindent (Dans SGA 4$^{1/2}$ : Th.\ Finitude : Theoremes de finitude en cohomologie $\ell$-adique.

Cycle : La classe de cohomologie associee a un cycle.

Rapport : Rapport sur la formule des traces).

\bigskip
\hrule
\bigskip

\part*{EXPOSE V}
\addcontentsline{toc}{part}{Expose V}

\begin{center}
\Large\bf SYSTEMES PROJECTIFS J-ADIQUES
\end{center}

\begin{center}
par J.P.\ JOUANOLOU
\end{center}

\tableofcontents

\newpage

\section{Generalites sur les A-categories abeliennes}

\subsection{Definitions et premieres proprietes}

\medskip
\noindent\textbf{1.1.1.} Soit $A$ un anneau commutatif. Une \underline{$A$-categorie} $\cC$ est une categorie munie, pour tout couple $(X,Y)$ d'objets de $\cC$, d'une structure de $A$-module sur l'ensemble des fleches $\Hom(X,Y)$ de sorte que, pour trois objets $X$, $Y$, $Z$ de $\cC$, la composition
\[
\Hom(Y,Z) \times \Hom(X,Y) \longrightarrow \Hom(X,Z)
\]
soit $A$-bilineaire. On dit que $\cC$ est \underline{additive} si elle est munie d'une structure de categorie additive dont le groupe abelien sous-jacent a la structure de $A$-module consideree soit le groupe additif naturel; elle est alors en particulier $\bZ$-lineaire.

\medskip
\noindent\textbf{1.1.2.} Une \underline{$A$-categorie abelienne} $\cC$ est une $A$-categorie additive dans laquelle tout morphisme admet un noyau et un conoyau, et dans laquelle tout morphisme se factorise canoniquement en un epimorphisme suivi d'un monomorphisme. Plus precisement, on suppose donnee une equivalence entre $\cC$ et une sous-categorie pleine de la categorie des faisceaux de $A$-modules sur un site, de sorte que les limites inductives et projectives finies, noyaux, conoyaux, images, coimages des morphismes, soient representables et se calculent ``faisceautiquement''.

\medskip
\noindent\textbf{1.1.3.} Un \underline{$A$-foncteur} (resp.\ \underline{$A$-foncteur additif}) entre deux $A$-categories (resp.\ $A$-categories additives) $F : \cC \longrightarrow \cC'$ est un foncteur (resp.\ foncteur additif) tel que, pour tout couple $(X,Y)$ d'objets de $\cC$, l'application $\Hom(X,Y) \longrightarrow \Hom(FX,FY)$ soit $A$-lineaire.

\medskip
\noindent\textbf{1.1.4.} Un \underline{$(A,B)$-foncteur} (ou $A$-$B$-foncteur) d'une $A$-categorie dans une $B$-categorie est la donnee d'un morphisme d'anneaux $A \longrightarrow B$ et d'un $B$-foncteur de $\cC$ dans $\cC'$; lorsque $A=B$ et que le morphisme $A \longrightarrow B$ est l'identite, on dit que c'est un $A$-foncteur.

\medskip
\noindent\textbf{1.1.5.} Soit $\cC$ une $A$-categorie abelienne. Un objet $M$ de $\cC$ est dit \underline{de type fini} (resp.\ \underline{noetherien}) si, pour tout systeme inductif filtrant $(N_i)_{i \in I}$ de sous-objets de $M$, le morphisme canonique
\[
\varinjlim_i \Hom(M,N_i) \longrightarrow \Hom(M, \varinjlim_i N_i)
\]
est un isomorphisme (resp.\ si tout sous-objet de $M$ est de type fini).

\medskip
\noindent\textbf{1.1.6.} Un objet $M$ de $\cC$ est dit \underline{coherent} s'il est de type fini et si, pour toute fleche $u : N \longrightarrow M$ dont la source est de type fini, le noyau de $u$ est de type fini. Lorsque $\cC$ est la categorie des modules sur un anneau, on retrouve les notions usuelles.

\medskip
\noindent\textbf{1.1.7.} Soient $\cC$ une $A$-categorie abelienne, $M$ un objet de $\cC$, et $J$ un ideal de $A$. On appelle \underline{$J$-filtration canonique} de $M$ la filtration decroissante definie par $J^n M$ ($n \in \bN$). On dit que $M$ est \underline{$J$-adique} si cette filtration est $J$-adique, i.e.\ si pour tout entier $n$, le morphisme canonique $J \otimes_A J^n M \longrightarrow J^{n+1} M$ est surjectif et si les morphismes $M/J^{n+1} M \longrightarrow M/J^n M$ sont surjectifs pour tout $n$.

\medskip
\noindent\textbf{1.1.8.} Dans toute la suite, nous ferons les conventions et hypotheses suivantes :

\textbf{(i)} $A$ designe un anneau commutatif noetherien.

\textbf{(ii) (AR)} $J$ est un ideal de $A$ tel que, pour tout $A$-module de type fini $M$, la $J$-filtration de $M$ soit $J$-adique (condition d'Artin-Rees).

\medskip
\noindent\textbf{1.1.9.} Soit $\cC$ une $A$-categorie abelienne. Un objet $M$ de $\cC$ est dit de \underline{type fini} (resp.\ \underline{noetherien}) s'il verifie la condition de 1.1.5. Lorsque $\cC$ est la categorie des $A$-modules, on retrouve les notions usuelles. On verifie immediatement que la sous-categorie pleine de $\cC$ formee des objets noetheriens est une sous-categorie abelienne epaisse (stable par sous-objets, quotients et extensions).

\subsection{Modules sur les anneaux filtres}

\medskip
\noindent\textbf{1.2.1.} Soit $A$ un anneau commutatif muni d'une filtration decroissante $F^\bullet A$, telle que $F^0 A = A$, et $F^p A \cdot F^q A \subset F^{p+q} A$ pour tout $p,q$. On note $\mathrm{gr}(A) = \bigoplus_{n \geq 0} F^n A/F^{n+1} A$ l'anneau gradue associe. On suppose que $A$ est separe et complet pour la filtration $F^\bullet$, et que $\mathrm{gr}(A)$ est noetherien.

\medskip
\noindent\textbf{1.2.2.} Soit $\cC$ une $A$-categorie abelienne. Un \underline{objet filtre} de $\cC$ est la donnee d'un objet $M$ de $\cC$ muni d'une filtration decroissante $F^\bullet M$, compatible avec la filtration de $A$, i.e.\ telle que $F^p A \cdot F^q M \subset F^{p+q} M$ pour tout $p,q$. On note $\mathrm{gr}(M) = \bigoplus_{n \geq 0} F^n M/F^{n+1} M$ le gradue associe, qui est un $\mathrm{gr}(A)$-module gradue.

\medskip
\noindent\textbf{1.2.3.} Soient $\cC$ une $A$-categorie abelienne, et $M$ un objet de $\cC$. On dit qu'un sous-objet $N$ de $M$ est \underline{purement de poids $n$} (relativement a la $J$-filtration) si $J^{n+1} N = 0$ et $N$ est un $A/J^{n+1}$-module fidele (au sens ou $N$ n'est pas annule par $J^n$).

\subsection{Categories de systemes projectifs}

\medskip
\noindent\textbf{1.3.1.} Soit $\cC$ une categorie. On note $\Hom(\bN^\circ, \cC)$ la categorie des systemes projectifs d'objets de $\cC$ indexes par l'ensemble ordonne $\bN$ des entiers naturels. Un objet de cette categorie est donc la donnee d'une suite $(M_n)_{n \in \bN}$ d'objets de $\cC$ et de morphismes de transition $u_{nm} : M_m \longrightarrow M_n$ pour $m \geq n$, verifiant les conditions de compatibilite usuelles.

\medskip
\noindent\textbf{1.3.2.} Soit $u : X \longrightarrow Y$ un morphisme de $\Hom(\bN^\circ, \cC)$. On dit que $u$ est un \underline{monomorphisme essentiel} (resp.\ \underline{epimorphisme essentiel}) si, pour tout entier $n$, il existe un entier $m \geq n$ tel que le morphisme $X_m \longrightarrow X_n \times_{Y_n} Y_m$ soit un monomorphisme (resp.\ un epimorphisme).

\medskip
\noindent\textbf{1.3.3.} On dit qu'un systeme projectif $X = (X_n)$ est \underline{essentiellement nul} si le morphisme identique $\mathrm{id}_X$ est a la fois un monomorphisme essentiel et un epimorphisme essentiel, ce qui equivaut a dire que pour tout entier $n$, il existe un entier $m \geq n$ tel que le morphisme de transition $u_{nm} : X_m \longrightarrow X_n$ soit nul.

\medskip
\noindent\textbf{1.3.4.} On verifie aisement que, lorsque $\cC$ est une categorie abelienne, les monomorphismes essentiels (resp.\ les epimorphismes essentiels) sont stables par composition et par produit fibre (resp.\ somme amalgamee). En particulier, la sous-categorie pleine de $\Hom(\bN^\circ, \cC)$ formee des systemes projectifs essentiellement nuls est une sous-categorie epaisse.

\medskip
\noindent\textbf{1.3.5.} On dit qu'un systeme projectif $X$ est \underline{AR-nul} s'il existe un entier $r$ tel que, pour tout $n$, le morphisme de transition $u_{n,n+r} : X_{n+r} \longrightarrow X_n$ soit nul. Il est clair qu'un systeme projectif AR-nul est essentiellement nul, et que tout sous-systeme projectif et tout quotient d'un systeme projectif AR-nul sont AR-nuls.

\medskip
\noindent\textbf{1.3.6.} \textbf{LEMME.} \underline{Si $u : X \longrightarrow Y$ est un morphisme de systemes projectifs essentiellement nuls, le noyau et le conoyau de $u$ sont essentiellement nuls.}

\textbf{Preuve} : Pour le noyau, cela resulte de la stabilite des monomorphismes essentiels par produit fibre. Pour le conoyau, cela resulte de la stabilite des epimorphismes essentiels par somme amalgamee.

\medskip
\noindent\textbf{1.3.7.} \textbf{LEMME.} \underline{Soit $0 \longrightarrow X' \longrightarrow X \longrightarrow X'' \longrightarrow 0$ une suite exacte de systemes projectifs. Alors :}
\begin{itemize}
\item[(i)] \underline{Si $X'$ et $X''$ sont essentiellement nuls, $X$ est essentiellement nul.}
\item[(ii)] \underline{Si $X$ est essentiellement nul, $X''$ est essentiellement nul.}
\item[(iii)] \underline{Si $X$ et $X''$ sont essentiellement nuls, $X'$ est essentiellement nul.}
\end{itemize}

\textbf{Preuve} : (i) resulte de la stabilite par composition. Pour (ii), il suffit de noter que si $X$ est essentiellement nul, alors pour tout $n$, il existe $m$ tel que $X_m \longrightarrow X_n$ soit nul, donc $X''_m \longrightarrow X''_n$ aussi. Pour (iii), on utilise la suite exacte des noyaux et le fait que $X'$ est essentiellement nul comme extension de systemes essentiellement nuls.

\medskip
\noindent\textbf{1.3.8.} \textbf{LEMME.} \underline{Soit $0 \longrightarrow X' \longrightarrow X \longrightarrow X'' \longrightarrow 0$ une suite exacte de systemes projectifs. Alors $X''$ est AR-nul si et seulement si $X$ est essentiellement nul et $X'$ est AR-nul.}

\textbf{Preuve} : Si $X''$ est AR-nul, il est essentiellement nul, donc d'apres 1.3.7 (ii), $X$ est essentiellement nul. De plus, pour tout $n$, il existe $m$ tel que $X_m \longrightarrow X_n$ soit nul, et comme $X''$ est AR-nul, il existe $r$ tel que $X''_{n+r} \longrightarrow X''_n$ soit nul; on en deduit que $X'$ est AR-nul. La reciproque est immediate.

\medskip
\noindent\textbf{1.3.9.} \textbf{LEMME.} \underline{Soit $0 \longrightarrow X' \longrightarrow X \longrightarrow X'' \longrightarrow 0$ une suite exacte de systemes projectifs telle que $X$ soit AR-nul et que les morphismes de transition de $X''$ soient des monomorphismes. Alors $X'$ est AR-nul.}

\textbf{Preuve} : Comme $X$ est AR-nul, il existe $r$ tel que $X_{n+r} \longrightarrow X_n$ soit nul pour tout $n$. Comme les morphismes de transition de $X''$ sont des monomorphismes, le morphisme $X'_{n+r} \longrightarrow X'_n$ est nul pour tout $n$, donc $X'$ est AR-nul.

\medskip
\noindent\textbf{1.3.10.} Soit $\cC$ une categorie abelienne. On note
\[
\Hom(\bN^\circ, \cC)_{AR\text{-}nul}
\]
la sous-categorie pleine de $\Hom(\bN^\circ, \cC)$ formee des systemes projectifs AR-nuls. D'apres 1.3.5, c'est une sous-categorie epaisse. On peut donc former la categorie quotient
\[
\Hom(\bN^\circ, \cC) / \Hom(\bN^\circ, \cC)_{AR\text{-}nul}\;.
\]
On appelle \underline{morphisme AR-isomorphisme} un morphisme de $\Hom(\bN^\circ, \cC)$ dont le noyau et le conoyau sont AR-nuls. Il resulte de la theorie generale des categories quotients qu'un morphisme de $\Hom(\bN^\circ, \cC)$ devient un isomorphisme dans la categorie quotient si et seulement si c'est un AR-isomorphisme.

\medskip
\noindent\textbf{1.3.11.} Soit $\cC$ une categorie abelienne. On dit qu'un objet $M$ de $\cC$ est \underline{filtrable} s'il admet une filtration finie dont les gradues sont purs (i.e.\ annules par une puissance de l'ideal $J$).

\medskip
\noindent\textbf{1.3.12.} \textbf{PROPOSITION.} \underline{Soit $\cC$ une $A$-categorie abelienne. Les systemes projectifs AR-nuls forment une sous-categorie localisante au sens de Gabriel.}

\textbf{Preuve} : Il s'agit de verifier que, pour tout systeme projectif $X$, il existe un sous-systeme projectif $Y$ qui soit maximal parmi les sous-systemes projectifs AR-nuls de $X$. On construit $Y_n$ comme la reunion des images des $X_m \longrightarrow X_n$ pour $m$ assez grand. Les details sont laisses au lecteur.

\medskip
\noindent\textbf{1.3.13.} On deduit de 1.3.12 que la categorie quotient $\Hom(\bN^\circ, \cC) / \Hom(\bN^\circ, \cC)_{AR\text{-}nul}$ est abelienne, et que le foncteur canonique
\[
Q : \Hom(\bN^\circ, \cC) \longrightarrow \Hom(\bN^\circ, \cC) / \Hom(\bN^\circ, \cC)_{AR\text{-}nul}
\]
est exact.

\medskip
\noindent\textbf{1.3.14.} Soient $\cC$ une $A$-categorie abelienne, et $\cC'$ une sous-categorie pleine de $\cC$. On note $\Hom_{AR}(\bN^\circ, \cC')$ la sous-categorie pleine de $\Hom(\bN^\circ, \cC)$ formee des systemes projectifs AR-isomorphes a des objets de $\Hom(\bN^\circ, \cC')$. Si $\cC'$ est stable par noyaux, conoyaux et extensions, il resulte de la theorie des categories quotients que $\Hom_{AR}(\bN^\circ, \cC')$ est une sous-categorie abelienne epaisse de $\Hom(\bN^\circ, \cC)$, et que l'on a une equivalence de categories
\[
\Hom_{AR}(\bN^\circ, \cC') / \Hom(\bN^\circ, \cC')_{AR\text{-}nul} \xrightarrow{\sim} \Hom(\bN^\circ, \cC') / \Hom(\bN^\circ, \cC')_{AR\text{-}nul}\;.
\]


\subsection{Exemples et constructions}

\medskip
\noindent\textbf{1.4.1.} Soit $M$ un objet d'une $A$-categorie abelienne $\cC$. On note $\kappa(M)$ (ou simplement $\kappa$) le systeme projectif constant de valeur $M$, i.e.\ le systeme projectif defini par $M_n = M$ pour tout $n$, les morphismes de transition etant l'identite de $M$. Le foncteur $\kappa : \cC \longrightarrow \Hom(\bN^\circ, \cC)$ est pleinement fidele et exact.

\medskip
\noindent\textbf{1.4.2.} Soit $r$ un entier $\geq 0$. On note $[r] : \Hom(\bN^\circ, \cC) \longrightarrow \Hom(\bN^\circ, \cC)$ le foncteur de "decalage" defini par $(X[r])_n = X_{n+r}$, avec les morphismes de transition induits par ceux de $X$. On verifie immediatement que $[r]$ est exact. De plus, on a un monomorphisme canonique $X[r] \longrightarrow X$, dont le conoyau est le systeme projectif $X_0 \longrightarrow X_1 \longrightarrow \cdots \longrightarrow X_{r-1} \longrightarrow 0 \longrightarrow 0 \longrightarrow \cdots$, qui est AR-nul.

\medskip
\noindent\textbf{1.4.3.} Soit $X = (X_n)$ un systeme projectif. On appelle \underline{limite projective} de $X$ (et on note $\varprojlim X$ ou $\varprojlim_n X_n$) l'objet de $\cC$ defini par la limite projective des $X_n$, lorsqu'elle existe. Si $\cC$ est complete, le foncteur $\varprojlim : \Hom(\bN^\circ, \cC) \longrightarrow \cC$ est un foncteur additif, exact a gauche. Si $R^q \varprojlim$ designent ses foncteurs derives a droite, on a, pour tout systeme projectif $X$, une suite exacte
\[
0 \longrightarrow R^1 \varprojlim_n X_{n-1} \longrightarrow \varprojlim_n X_n / X_{n+1} \longrightarrow \varprojlim_n X_n / X_{n+1} \longrightarrow R^2 \varprojlim_n X_{n-1} \longrightarrow 0
\]
et des isomorphismes $R^q \varprojlim_n X_{n-1} \xrightarrow{\sim} R^{q+2} \varprojlim_n X_{n-1}$ pour $q \geq 1$.

\medskip
\noindent\textbf{1.4.4.} Soient $\cC$ une $A$-categorie abelienne, et $J$ un ideal de $A$. Pour tout objet $M$ de $\cC$, on note $J^n M$ ($n \geq 0$) le sous-objet image du morphisme canonique $J^n \otimes_A M \longrightarrow M$. On appelle \underline{$J$-filtration canonique} de $M$ la filtration decroissante definie par les $J^n M$. Le \underline{separe complete $J$-adique} de $M$ est la limite projective $\varprojlim_n M/J^n M$, lorsqu'elle existe.

\medskip
\noindent\textbf{1.4.5.} Soit $\hat{\cC}$ la categorie des objets $J$-adiquement complets et separes de $\cC$. On suppose que $\cC$ est complete et que les produits denombrables sont exacts. Alors le foncteur de completion $M \mapsto \hat{M} = \varprojlim_n M/J^n M$ est exact. De plus, pour tout objet $M$ de $\cC$, le morphisme canonique $\hat{M} \longrightarrow \varprojlim_n M/J^n M$ est un isomorphisme.

\medskip
\noindent\textbf{1.4.6.} On suppose que $\cC$ est la categorie des $A$-modules, avec $A$ noetherien, et que $J$ verifie la condition (AR) de 1.1.8 (ii). Soit $M$ un $A$-module de type fini. Alors :
\begin{itemize}
\item[(i)] La $J$-filtration de $M$ est $J$-adique.
\item[(ii)] Le separe complete $\hat{M} = \varprojlim_n M/J^n M$ est un $A$-module de type fini.
\item[(iii)] Le morphisme canonique $\hat{A} \otimes_A M \longrightarrow \hat{M}$ est un isomorphisme.
\end{itemize}

\medskip
\noindent\textbf{1.4.7.} Sous les hypotheses de 1.4.6, soit $M$ un $A$-module de type fini, et $N$ un sous-module de $M$. Alors la filtration induite sur $N$ par la $J$-filtration de $M$ est $J$-adique. Autrement dit, pour tout entier $n$, il existe un entier $m$ tel que $N \cap J^m M \subset J^n N$.

\medskip
\noindent\textbf{1.4.8.} \textbf{LEMME.} \underline{Soit $M$ un objet d'une $A$-categorie abelienne $\cC$, et $(N_i)_{i \in I}$ un systeme inductif filtrant de sous-objets de $M$. Alors le morphisme canonique}
\[
\varinjlim_i \Hom(M,N_i) \longrightarrow \Hom(M, \varinjlim_i N_i)
\]
\underline{est un isomorphisme si et seulement si, pour tout morphisme $u : M \longrightarrow \varinjlim_i N_i$, il existe un indice $i$ et un morphisme $v : M \longrightarrow N_i$ relevant $u$.}

\textbf{Preuve} : C'est une reformulation de la definition.

\medskip
\noindent\textbf{1.4.9.} \textbf{LEMME.} \underline{Soit $M$ un objet noetherien d'une $A$-categorie abelienne $\cC$. Alors, pour tout systeme inductif filtrant $(N_i)_{i \in I}$ de sous-objets de $M$, le morphisme canonique}
\[
\varinjlim_i \Hom(M, N_i) \longrightarrow \Hom(M, \varinjlim_i N_i)
\]
\underline{est un isomorphisme.}

\textbf{Preuve} : Soit $u : M \longrightarrow \varinjlim_i N_i$. Comme $M$ est noetherien, l'image de $u$ est de type fini, donc est contenue dans un $N_i$. Il existe donc un relevement $v : M \longrightarrow N_i$.

\medskip
\noindent\textbf{1.4.10.} \textbf{LEMME.} \underline{Soit $M$ un objet noetherien d'une $A$-categorie abelienne $\cC$. Alors tout sous-objet de $M$ est noetherien.}

\textbf{Preuve} : Soit $N$ un sous-objet de $M$, et $(P_i)_{i \in I}$ un systeme inductif filtrant de sous-objets de $N$. C'est aussi un systeme inductif filtrant de sous-objets de $M$. Comme $M$ est noetherien, le morphisme $\varinjlim_i \Hom(M,P_i) \longrightarrow \Hom(M, \varinjlim_i P_i)$ est un isomorphisme. Par restriction, on en deduit que $\varinjlim_i \Hom(N,P_i) \longrightarrow \Hom(N, \varinjlim_i P_i)$ est un isomorphisme.

\medskip
\noindent\textbf{1.4.11.} On appelle \underline{$J$-adique} (resp.\ \underline{AR-$J$-adique}) un systeme projectif $X = (X_n)$ tel que, pour tout $n$, $X_n$ soit annule par $J^{n+1}$ et que les morphismes de transition $X_m \longrightarrow X_n$ pour $m \geq n$ induisent des isomorphismes $X_m / J^{n+1} X_m \xrightarrow{\sim} X_n$. Un tel systeme projectif est dit \underline{strict} si les morphismes de transition sont surjectifs.

\medskip
\noindent\textbf{1.4.12.} Soit $X = (X_n)$ un systeme projectif $J$-adique. On dit que $X$ est \underline{noetherien} (resp.\ \underline{artinien}) si les $X_n$ sont des objets noetheriens (resp.\ artiniens) de $\cC$. On notera que, si $X$ est noetherien, alors $X_0$ est noetherien, et, pour tout $n$, le morphisme $X_n \longrightarrow X_0$ est surjectif.

\medskip
\noindent\textbf{1.4.13.} On appelle \underline{topologie $J$-adique} sur un objet $M$ d'une $A$-categorie abelienne $\cC$ la topologie definie par les sous-objets $J^n M$ ($n \geq 0$). On dit que $M$ est \underline{separe} (resp.\ \underline{complet}) pour cette topologie si le morphisme canonique $M \longrightarrow \varprojlim_n M/J^n M$ est un monomorphisme (resp.\ un epimorphisme).

\medskip
\noindent\textbf{1.4.14.} \textbf{LEMME.} \underline{Soit $M$ un objet d'une $A$-categorie abelienne $\cC$. On suppose que $M$ est separe et complet pour la topologie $J$-adique. Alors le systeme projectif $(M/J^n M)_n$ est strict $J$-adique, et son image dans la categorie quotient par les systemes AR-nuls est isomorphe a l'image du systeme constant $\kappa(M)$.}

\textbf{Preuve} : Comme $M$ est separe et complet, le morphisme canonique $M \longrightarrow \varprojlim_n M/J^n M$ est un isomorphisme. Les morphismes de transition $M/J^m M \longrightarrow M/J^n M$ pour $m \geq n$ sont surjectifs, et leur noyau est $J^n M / J^m M$. La condition $J$-adique resulte de la definition.

\section{Condition de Mittag-Leffler-Artin-Rees et calcul de fractions}

\medskip
\noindent\textbf{2.0.} Dans cette section, $\cC$ designe une $A$-categorie abelienne, et $J$ un ideal de $A$ verifiant la condition (AR) de 1.1.8 (ii).

\subsection{Condition de Mittag-Leffler-Artin-Rees}

\medskip
\noindent\textbf{2.1.1.} Soit $X = (X_n)_{n \in \bN}$ un systeme projectif d'objets de $\cC$. On dit que $X$ \underline{verifie la condition de Mittag-Leffler-Artin-Rees} (en abrege \underline{(MLAR)}) si, pour tout entier $n$, il existe un entier $m \geq n$ tel que, pour tout entier $p \geq m$, l'image de $X_p \longrightarrow X_n$ soit egale a l'image de $X_m \longrightarrow X_n$. Autrement dit, la suite des images $\mathrm{Im}(X_m \longrightarrow X_n)_{m \geq n}$ est stationnaire.

\medskip
\noindent\textbf{2.1.2.} On dit que $X$ \underline{verifie la condition de Mittag-Leffler-Artin-Rees stricte} (en abrege \underline{(MLAR stricte)}) si, pour tout entier $n$, il existe un entier $m \geq n$ tel que, pour tout entier $p \geq m$, l'image de $X_p \longrightarrow X_n$ soit egale a l'image de $X_m \longrightarrow X_n$, et de plus cette image soit un objet $J$-adique de $\cC$.

\medskip
\noindent\textbf{2.1.3.} \textbf{LEMME.} \underline{Soit $X$ un systeme projectif verifiant (MLAR). Alors il existe un systeme projectif strict $Y$, AR-isomorphe a $X$, tel que les morphismes de transition de $Y$ soient des epimorphismes.}

\textbf{Preuve} : Pour tout $n$, soit $m(n)$ un entier tel que $\mathrm{Im}(X_m \longrightarrow X_n)$ soit constante pour $m \geq m(n)$. Posons $Y_n = \mathrm{Im}(X_{m(n)} \longrightarrow X_n)$. Les morphismes de transition $Y_{n+1} \longrightarrow Y_n$ sont surjectifs par construction, et le morphisme canonique $Y \longrightarrow X$ est un AR-isomorphisme.

\medskip
\noindent\textbf{2.1.4.} \textbf{LEMME.} \underline{Soit $0 \longrightarrow X' \longrightarrow X \longrightarrow X'' \longrightarrow 0$ une suite exacte de systemes projectifs.}
\begin{itemize}
\item[(i)] \underline{Si $X$ verifie (MLAR), alors $X''$ verifie (MLAR).}
\item[(ii)] \underline{Si $X'$ et $X''$ verifient (MLAR), alors $X$ verifie (MLAR).}
\end{itemize}

\textbf{Preuve} : (i) est immediat, car les images dans $X''$ sont les images des images dans $X$. Pour (ii), cela resulte du lemme des cinq applique aux suites exactes des images.

\medskip
\noindent\textbf{2.1.5.} \textbf{LEMME.} \underline{Soit $X$ un systeme projectif $J$-adique. Pour que $X$ verifie (MLAR), il faut et il suffit que $X$ soit strict.}

\textbf{Preuve} : Si $X$ est strict, les morphismes de transition sont surjectifs, donc (MLAR) est verifiee trivialement. Reciproquement, si $X$ verifie (MLAR), pour tout $n$, il existe $m$ tel que $\mathrm{Im}(X_p \longrightarrow X_n) = \mathrm{Im}(X_m \longrightarrow X_n)$ pour tout $p \geq m$. Comme $X$ est $J$-adique, cela implique que les morphismes de transition sont surjectifs.

\medskip
\noindent\textbf{2.1.6.} \textbf{LEMME.} \underline{Soit $X$ un systeme projectif noetherien. Alors $X$ verifie (MLAR).}

\textbf{Preuve} : Pour tout $n$, les images $\mathrm{Im}(X_m \longrightarrow X_n)$ forment une suite croissante de sous-objets de $X_n$. Comme $X_n$ est noetherien, cette suite est stationnaire.

\subsection{Systemes projectifs stricts}

\medskip
\noindent\textbf{2.2.1.} Soit $X = (X_n)$ un systeme projectif strict $J$-adique. On appelle \underline{noyau} de $X$ le systeme projectif $\Ker(X) = (\Ker(X_n \longrightarrow X_0))_n$. C'est un sous-systeme projectif strict de $X$, et le quotient $X / \Ker(X)$ est le systeme projectif constant $\kappa(X_0)$.

\medskip
\noindent\textbf{2.2.2.} Soient $X = (X_n)$ un systeme projectif strict $J$-adique, et $Y = (Y_n)$ un sous-systeme projectif strict. On dit que $Y$ est \underline{$J$-adapte} a $X$ si, pour tout entier $n$, on a $Y_n = X_n \cap Y_0$ dans $X_0$ (ou $X_n$ est considere comme sous-objet de $X_0$ via le morphisme de transition $X_n \longrightarrow X_0$).

\medskip
\noindent\textbf{2.2.3.} \textbf{LEMME.} \underline{Soient $X$ un systeme projectif strict $J$-adique, et $Y$ un sous-systeme projectif strict $J$-adapte. Alors $Y$ est $J$-adique si et seulement si, pour tout entier $n$, le morphisme canonique $J^n \otimes_A Y_0 \longrightarrow Y_0 \cap J^n X_0$ est surjectif.}

\textbf{Preuve} : C'est une reformulation des definitions.

\medskip
\noindent\textbf{2.2.4.} \textbf{LEMME.} \underline{Soit $X$ un systeme projectif strict $J$-adique, et $Y$ un sous-systeme projectif strict $J$-adapte. Alors $Y$ est $J$-adique si et seulement si le gradue $
\mathrm{gr}(Y) = \bigoplus_{n \geq 0} Y_n/Y_{n+1}$
est un $\mathrm{gr}_J(A)$-module gradue.}

\textbf{Preuve} : Par definition, $\mathrm{gr}(Y)$ est un $\mathrm{gr}_J(A)$-module gradue si et seulement si l'action de $J^n/J^{n+1}$ sur $Y_m/Y_{m+1}$ se factorise par $Y_{n+m}/Y_{n+m+1}$, ce qui equivaut a la condition de 2.2.3.

\medskip
\noindent\textbf{2.2.5.} Soit $X = (X_n)$ un systeme projectif strict $J$-adique. On appelle \underline{gradue associe} a $X$ le $\mathrm{gr}_J(A)$-module gradue
\[
\mathrm{gr}(X) = \bigoplus_{n \geq 0} X_n / X_{n+1}\;,
\]
où $X_n / X_{n+1}$ est de degre $n$. Le produit $J^p/J^{p+1} \times X_n/X_{n+1} \longrightarrow X_{n+p}/X_{n+p+1}$ est induit par la structure de $A$-module de $X_0$.

\medskip
\noindent\textbf{2.2.6.} \textbf{LEMME.} \underline{Soit $X$ un systeme projectif strict $J$-adique. Alors $X$ est noetherien (resp.\ artinien) si et seulement si son gradue $\mathrm{gr}(X)$ est un $\mathrm{gr}_J(A)$-module noetherien (resp.\ artinien).}

\textbf{Preuve} : Supposons $\mathrm{gr}(X)$ noetherien. Soit $Y$ un sous-systeme projectif strict de $X$. Alors $\mathrm{gr}(Y)$ est un sous-module gradue de $\mathrm{gr}(X)$, donc est de type fini. Il en resulte que $Y$ est noetherien. La reciproque est immediate.

\medskip
\noindent\textbf{2.2.7.} \textbf{LEMME.} \underline{Soit $X$ un systeme projectif strict $J$-adique. Alors, pour tout entier $r \geq 0$, le systeme projectif $X[r]$ est AR-isomorphe a $X$.}

\textbf{Preuve} : Le morphisme canonique $X[r] \longrightarrow X$ est un monomorphisme, et son conoyau est AR-nul car les composantes de rang $< r$ sont nulles.

\subsection{Calcul de fractions. Sous-categories localisantes}

\medskip
\noindent\textbf{2.3.1.} Soit $\cC$ une categorie abelienne. On appelle \underline{sous-categorie localisante} (au sens de Gabriel) de $\cC$ une sous-categorie pleine $\cL$ de $\cC$ verifiant les conditions suivantes :
\begin{itemize}
\item[(i)] $\cL$ est une sous-categorie abelienne epaisse (stable par sous-objets, quotients et extensions).
\item[(ii)] $\cL$ est stable par sous-quotients.
\item[(iii)] Tout objet de $\cC$ contient un sous-objet maximal appartenant a $\cL$.
\end{itemize}

\medskip
\noindent\textbf{2.3.2.} Soit $\cL$ une sous-categorie localisante de $\cC$. On dit qu'un morphisme $f : X \longrightarrow Y$ de $\cC$ est une \underline{$\cL$-equivalence} si son noyau et son conoyau appartiennent a $\cL$. On verifie que la classe des $\cL$-equivalences admet un calcul de fractions a droite et a gauche.

\medskip
\noindent\textbf{2.3.3.} La categorie quotient $\cC/\cL$ est definie comme la localisation de $\cC$ par rapport aux $\cL$-equivalences. On a un foncteur canonique $Q : \cC \longrightarrow \cC/\cL$ qui est exact et universal pour les foncteurs exacts envoyant les objets de $\cL$ sur zero.

\medskip
\noindent\textbf{2.3.4.} \textbf{PROPOSITION.} \underline{Soient $\cC$ une $A$-categorie abelienne, et $\cE$ une sous-categorie localisante de $\Hom(\bN^\circ, \cC)$ formee de systemes projectifs essentiellement nuls. Alors la categorie quotient $\Hom(\bN^\circ, \cC) / \cE$ est abelienne, et le foncteur canonique est exact.}

\textbf{Preuve} : C'est un cas particulier de la theorie de Gabriel.

\medskip
\noindent\textbf{2.3.5.} Soient $\cC$ une $A$-categorie abelienne, et $\cC'$ une sous-categorie pleine de $\cC$. On suppose que $\cC'$ est une sous-categorie abelienne epaisse de $\cC$. On note $\Hom_{AR}(\bN^\circ, \cC')$ la sous-categorie pleine de $\Hom(\bN^\circ, \cC)$ formee des systemes projectifs AR-isomorphes a des objets de $\Hom(\bN^\circ, \cC')$. On verifie que $\Hom_{AR}(\bN^\circ, \cC')$ est une sous-categorie localisante de $\Hom(\bN^\circ, \cC)$.

\medskip
\noindent\textbf{2.3.6.} \textbf{LEMME.} \underline{Soit $\cC$ une $A$-categorie abelienne, et $\cE$ la sous-categorie localisante des systemes essentiellement nuls. Alors, pour tout objet $X$ de $\Hom(\bN^\circ, \cC)$, le morphisme canonique $X \longrightarrow Q(X)$ est essentiellement surjectif au sens suivant : pour tout sous-objet $Y$ de $X$ tel que $Q(Y) = Q(X)$, il existe un sous-objet $Z$ de $Y$ tel que $X/Z$ soit essentiellement nul.}

\textbf{Preuve} : C'est la description du noyau de $Q$.

\medskip
\noindent\textbf{2.3.7.} \textbf{LEMME.} \underline{Soient $\cC$ une $A$-categorie abelienne, et $\cL$ une sous-categorie localisante de systemes essentiellement nuls. On suppose que $\cL$ contient les systemes AR-nuls. Alors, pour tout systeme projectif strict $J$-adique $X$, et tout sous-systeme projectif strict $Y$ de $X$, il existe un sous-systeme projectif strict $Z$ de $Y$, $J$-adapte, tel que $Y/Z \in \cL$.}

\textbf{Preuve} : On construit $Z$ comme l'intersection des noyaux des morphismes $Y_n \longrightarrow X_0/J^{n+1}X_0$. Les details sont laisses au lecteur.

\medskip
\noindent\textbf{2.3.8.} \textbf{LEMME.} \underline{Soient $\cC$ une $A$-categorie abelienne, et $\cL$ la sous-categorie localisante des systemes essentiellement nuls. Alors, pour tout systeme projectif strict $J$-adique $X$, le foncteur $Y \mapsto Q(Y)$ de la categorie des sous-systemes projectifs stricts $J$-adaptes de $X$ dans la categorie des sous-objets de $Q(X)$ est une equivalence de categories.}

\textbf{Preuve} : La pleine fidelite resulte de ce que deux sous-systemes $J$-adaptes ayant meme image dans $Q(X)$ different par un systeme essentiellement nul. L'essentielle surjectivite resulte de 2.3.7.


\section{Systemes projectifs J-adiques et AR-J-adiques}

\medskip
\noindent\textbf{3.0.} Dans cette section, $\cC$ designe une $A$-categorie abelienne, et $J$ un ideal de $A$ verifiant la condition (AR) de 1.1.8 (ii).

\subsection{Definitions}

\medskip
\noindent\textbf{3.1.1.} On dit qu'un objet $X$ de $\Hom(\bN^\circ, \cC)$ est \underline{$J$-adique} s'il est isomorphe a un systeme projectif strict $J$-adique. On notera $J\text{-}\mathrm{ad}(\cC)$ (resp.\ $J\text{-}\mathrm{adf}(\cC)$) la sous-categorie pleine de $\Hom(\bN^\circ, \cC)$ formee des systemes projectifs $J$-adiques noetheriens (resp.\ artiniens).

\medskip
\noindent\textbf{3.1.2.} On dit qu'un objet $X$ de $\Hom(\bN^\circ, \cC)$ est \underline{AR-$J$-adique} s'il est AR-isomorphe a un systeme projectif strict $J$-adique. On notera $\mathrm{AR\text{-}J\text{-}ad}(\cC)$ (resp.\ $\mathrm{AR\text{-}J\text{-}adf}(\cC)$) la sous-categorie pleine de $\Hom(\bN^\circ, \cC)$ formee des systemes projectifs AR-$J$-adiques noetheriens (resp.\ artiniens).

\medskip
\noindent\textbf{3.1.3.} \textbf{LEMME.} \underline{Soit $X$ un systeme projectif $J$-adique. Les conditions suivantes sont equivalentes :}
\begin{itemize}
\item[(i)] \underline{$X$ est noetherien (resp.\ artinien).}
\item[(ii)] \underline{Le gradue $\mathrm{gr}(X)$ est un $\mathrm{gr}_J(A)$-module noetherien (resp.\ artinien).}
\item[(iii)] \underline{$X_0$ est un objet noetherien (resp.\ artinien) de $\cC$, et le $\mathrm{gr}_J(A)$-module gradue $\mathrm{gr}(X)$ est de type fini.}
\end{itemize}

\textbf{Preuve} : (i) $\Rightarrow$ (ii) resulte de 2.2.6. (ii) $\Rightarrow$ (iii) est immediat, car $X_0 = \mathrm{gr}_0(X)$. (iii) $\Rightarrow$ (i) : si $\mathrm{gr}(X)$ est de type fini, alors, pour tout sous-systeme projectif strict $Y$, le gradue $\mathrm{gr}(Y)$ est un sous-module gradue de $\mathrm{gr}(X)$, donc est de type fini, et par suite $Y$ est noetherien.

\medskip
\noindent\textbf{3.1.4.} \textbf{PROPOSITION.} \underline{Soit $X$ un systeme projectif AR-$J$-adique. Alors $X$ est AR-isomorphe a un systeme projectif strict $J$-adique dont les morphismes de transition sont des epimorphismes.}

\textbf{Preuve} : Par hypothese, il existe un systeme projectif strict $J$-adique $Y$ et un AR-isomorphisme $u : Y \longrightarrow X$. Le noyau et le conoyau de $u$ sont AR-nuls. Quitte a remplacer $Y$ par $Y/\Ker(u)$, on peut supposer $u$ injectif. Le conoyau etant AR-nul, on en deduit que $X$ est AR-isomorphe a $Y$.

\medskip
\noindent\textbf{3.1.5.} \textbf{LEMME.} \underline{Soit $X$ un systeme projectif $J$-adique noetherien. Alors, pour tout entier $r \geq 0$, le systeme projectif $X[r]$ est $J$-adique noetherien.}

\textbf{Preuve} : Comme $X$ est noetherien, les morphismes de transition sont surjectifs (lemme 2.1.6 et 2.1.5). Donc $X[r]$ est strict, et comme $(X[r])_n = X_{n+r}$, il est $J$-adique. De plus, $\mathrm{gr}(X[r])$ est isomorphe a $\mathrm{gr}(X)$ decale, donc est noetherien.

\medskip
\noindent\textbf{3.1.6.} \textbf{PROPOSITION.} \underline{Soit $X$ un systeme projectif AR-$J$-adique noetherien. Alors il existe un entier $r$ et un systeme projectif $J$-adique noetherien $Y$ tels que $X$ soit isomorphe a $Y[r]$ dans la categorie quotient par les systemes AR-nuls.}

\textbf{Preuve} : Par definition, il existe un systeme projectif $J$-adique noetherien $Z$ et un AR-isomorphisme $u : Z \longrightarrow X$. Comme $\Ker(u)$ et $\Coker(u)$ sont AR-nuls, il existe $r$ tel que $Z[r] \longrightarrow X[r]$ soit un isomorphisme. On prend $Y = Z[r]$.

\subsection{Le foncteur limite projective}

\medskip
\noindent\textbf{3.2.1.} Soit $X = (X_n)$ un systeme projectif d'objets d'une $A$-categorie abelienne $\cC$. On suppose que $\cC$ est complete et que les produits denombrables sont exacts. Alors le foncteur $\varprojlim$ est exact a gauche, et ses foncteurs derives $R^q \varprojlim$ ($q \geq 1$) verifient :
\begin{itemize}
\item[(i)] $R^q \varprojlim_n X_n = 0$ pour $q \geq 2$.
\item[(ii)] On a une suite exacte
\[
0 \longrightarrow R^1 \varprojlim_n X_n \longrightarrow \prod_n X_n \xrightarrow{\;1-u\;} \prod_n X_n \longrightarrow \varprojlim_n^{(1)} X_n \longrightarrow 0\;,
\]
où $u$ est le morphisme de decalage.
\end{itemize}

\medskip
\noindent\textbf{3.2.2.} \textbf{LEMME.} \underline{Soit $X$ un systeme projectif AR-nul. Alors $R^q \varprojlim X = 0$ pour tout $q \geq 0$.}

\textbf{Preuve} : Comme $X$ est AR-nul, il existe $r$ tel que $X_{n+r} \longrightarrow X_n$ soit nul pour tout $n$. Il en resulte que $R^q \varprojlim X = 0$ pour tout $q$.

\medskip
\noindent\textbf{3.2.3.} \textbf{PROPOSITION.} \underline{Soit $X$ un systeme projectif AR-$J$-adique noetherien. Alors $R^1 \varprojlim X = 0$.}

\textbf{Preuve} : D'apres 3.1.6, il existe $r$ et un systeme projectif $J$-adique noetherien $Y$ tels que $X$ soit AR-isomorphe a $Y[r]$. D'apres 3.2.2, $R^1 \varprojlim X \simeq R^1 \varprojlim Y[r]$. Or $R^1 \varprojlim Y[r] \simeq R^1 \varprojlim Y$ (par decalage), et ce dernier est nul car $Y$ est noetherien et verifie donc (MLAR) d'apres 2.1.6.

\medskip
\noindent\textbf{3.2.4.} \textbf{COROLLAIRE.} \underline{Soit $X$ un systeme projectif AR-$J$-adique noetherien. Alors $\varprojlim X$ est un objet $J$-adiquement complet et separe de $\cC$, et le morphisme canonique $X_n \longrightarrow \varprojlim X / J^{n+1} \varprojlim X$ est un isomorphisme pour tout $n$.}

\textbf{Preuve} : D'apres 3.2.3, la suite $0 \longrightarrow \varprojlim X \longrightarrow \prod_n X_n \longrightarrow \prod_n X_n \longrightarrow 0$ est exacte. Il en resulte que $\varprojlim X$ est complet et separe pour la topologie $J$-adique. La deuxieme assertion resulte de ce que $X$ est $J$-adique.

\subsection{Completion $J$-adique}

\medskip
\noindent\textbf{3.3.1.} Soit $\cC$ une $A$-categorie abelienne. On suppose que $\cC$ est complete et que les produits denombrables sont exacts. On appelle \underline{completion $J$-adique} le foncteur
\[
\Lambda_J : \cC \longrightarrow \cC, \quad M \longmapsto \varprojlim_n M/J^n M\;.
\]
Ce foncteur est additif et exact a gauche. On note $L_i \Lambda_J$ ($i \geq 0$) ses foncteurs derives a gauche.

\medskip
\noindent\textbf{3.3.2.} \textbf{LEMME.} \underline{Pour tout objet $M$ de $\cC$, on a $L_i \Lambda_J(M) = 0$ pour $i \geq 2$.}

\textbf{Preuve} : Les foncteurs derives de la limite projective sont nuls en degre $\geq 2$ (d'apres 3.2.1), et la completion est un compose de foncteurs $\varprojlim$ et de passage au quotient.

\medskip
\noindent\textbf{3.3.3.} \textbf{PROPOSITION.} \underline{Soit $M$ un objet noetherien de $\cC$. Alors $L_1 \Lambda_J(M) = 0$, et le morphisme canonique $\Lambda_J(M) \longrightarrow \varprojlim_n M/J^n M$ est un isomorphisme.}

\textbf{Preuve} : Comme $M$ est noetherien, la $J$-filtration de $M$ est $J$-adique (condition AR). Il en resulte que le systeme projectif $(M/J^n M)$ verifie (MLAR), donc $R^1 \varprojlim M/J^n M = 0$. On en deduit que $L_1 \Lambda_J(M) = 0$.

\medskip
\noindent\textbf{3.3.4.} \textbf{THEOREME.} \underline{Soit $\cC$ une $A$-categorie abelienne. On suppose que $\cC$ est complete et que les produits denombrables sont exacts. Soit $M$ un objet noetherien de $\cC$. Alors le morphisme canonique}
\[
\hat{A} \otimes_A M \longrightarrow \Lambda_J(M)
\]
\underline{est un isomorphisme.}

\textbf{Preuve} : On se ramene a supposer $A$ complet. Comme $M$ est noetherien, il existe une presentation finie $A^p \longrightarrow A^q \longrightarrow M \longrightarrow 0$. Par exactitude a droite du produit tensoriel, on obtient une suite exacte $\hat{A}^p \longrightarrow \hat{A}^q \longrightarrow \hat{A} \otimes_A M \longrightarrow 0$. Comme $A$ est noetherien, le foncteur de completion est exact sur les modules de type fini, donc on a une suite exacte $\hat{A}^p \longrightarrow \hat{A}^q \longrightarrow \hat{M} \longrightarrow 0$. On en deduit un isomorphisme $\hat{A} \otimes_A M \xrightarrow{\sim} \hat{M}$.

\section{Filtrations et graduations}

\medskip
\noindent\textbf{4.0.} Dans cette section, $\cC$ designe une $A$-categorie abelienne, et $J$ un ideal de $A$ verifiant la condition (AR) de 1.1.8 (ii).

\subsection{Filtrations et graduations sur les systemes projectifs}

\medskip
\noindent\textbf{4.1.1.} Soit $X = (X_n)$ un systeme projectif d'objets de $\cC$. Une \underline{filtration} de $X$ est la donnee, pour tout entier $n$, d'une filtration decroissante $F^\bullet X_n$ de $X_n$, de sorte que les morphismes de transition $X_m \longrightarrow X_n$ soient compatibles aux filtrations. On dit que la filtration est \underline{$J$-bonne} si, pour tout $n$, la filtration de $X_n$ est $J$-bonne, i.e.\ il existe $n_0$ tel que $J^p F^{n_0} X_n = F^{n_0+p} X_n$ pour tout $p \geq 0$.

\medskip
\noindent\textbf{4.1.2.} Soit $X = (X_n)$ un systeme projectif filtre. On appelle \underline{gradue associe} le systeme projectif $\mathrm{gr}(X) = (\mathrm{gr}(X_n))_n$, ou $\mathrm{gr}(X_n) = \bigoplus_{p \geq 0} F^p X_n / F^{p+1} X_n$. C'est un systeme projectif de $\mathrm{gr}_J(A)$-modules gradues.

\medskip
\noindent\textbf{4.1.3.} \textbf{LEMME.} \underline{Soit $X$ un systeme projectif strict $J$-adique muni d'une filtration $J$-bonne. Alors le gradue $\mathrm{gr}(X)$ est un systeme projectif strict $\mathrm{gr}_J(A)$-adique.}

\textbf{Preuve} : Comme la filtration est $J$-bonne, on a, pour $p$ assez grand, $F^{p+1} X_n = J \cdot F^p X_n$. Il en resulte que les morphismes de transition de $\mathrm{gr}(X)$ sont surjectifs et que $\mathrm{gr}(X)$ est $\mathrm{gr}_J(A)$-adique.

\medskip
\noindent\textbf{4.1.4.} \textbf{LEMME.} \underline{Soit $X$ un systeme projectif strict $J$-adique noetherien, muni d'une filtration $J$-bonne. Alors le gradue $\mathrm{gr}(X)$ est un systeme projectif strict $\mathrm{gr}_J(A)$-adique noetherien.}

\textbf{Preuve} : Comme $X$ est noetherien et la filtration $J$-bonne, le gradue $\mathrm{gr}(X_0)$ est un $\mathrm{gr}_J(A)$-module noetherien. D'apres 4.1.3, $\mathrm{gr}(X)$ est strict $\mathrm{gr}_J(A)$-adique. Comme $\mathrm{gr}(X_0)$ est noetherien, $\mathrm{gr}(X)$ est noetherien.

\subsection{Le lemme d'Artin-Rees}

\medskip
\noindent\textbf{4.2.1.} \textbf{(Lemme d'Artin-Rees).} \underline{Soit $M$ un objet noetherien d'une $A$-categorie abelienne $\cC$, et $N$ un sous-objet de $M$. On munit $M$ de sa $J$-filtration canonique. Alors la filtration induite sur $N$ est $J$-adique, i.e.\ pour tout entier $n$, il existe un entier $m$ tel que $N \cap J^m M \subset J^n N$.}

\textbf{Preuve} : La sous-categorie des objets noetheriens etant abelienne, $N$ est noetherien. Considerons le systeme projectif $(N / N \cap J^n M)_n$. C'est un sous-systeme projectif du systeme constant $\kappa(N)$. Comme $N$ est noetherien, ce systeme projectif verifie (MLAR) d'apres 2.1.6. Il existe donc, pour tout $n$, un entier $m$ tel que $N \cap J^m M \subset J^n N$.

\medskip
\noindent\textbf{4.2.2.} \textbf{LEMME.} \underline{Soit $M$ un objet noetherien de $\cC$. Alors, pour tout entier $n$, le morphisme canonique}
\[
J^n \otimes_A M \longrightarrow J^n M
\]
\underline{est surjectif.}

\textbf{Preuve} : Comme $M$ est noetherien, il est de type fini. La surjectivite resulte alors de la definition de $J^n M$ comme image du morphisme canonique.

\medskip
\noindent\textbf{4.2.3.} \textbf{LEMME.} \underline{Soit $M$ un objet noetherien de $\cC$. Alors le systeme projectif $(J^n M)_n$ est AR-nul si et seulement si $M$ est annule par une puissance de $J$.}

\textbf{Preuve} : Si $M$ est annule par $J^r$, alors $J^n M = 0$ pour $n \geq r$, donc le systeme projectif est AR-nul. Reciproquement, si $(J^n M)$ est AR-nul, il existe $r$ tel que $J^{n+r} M \longrightarrow J^n M$ soit nul pour tout $n$. Comme $M$ est noetherien, il existe $n$ tel que $J^n M = J^{n+1} M$. On en deduit que $J^n M = 0$.

\medskip
\noindent\textbf{4.2.4.} \textbf{PROPOSITION.} \underline{Soit $X$ un systeme projectif AR-$J$-adique noetherien, muni d'une filtration $J$-bonne. Alors le gradue $\mathrm{gr}(X)$ est AR-}$\mathrm{gr}_J(A)$-\underline{adique noetherien.}

\textbf{Preuve} : D'apres 3.1.6, on peut supposer que $X$ est $J$-adique noetherien. La proposition resulte alors de 4.1.4.

\medskip
\noindent\textbf{4.2.5.} \textbf{LEMME.} \underline{Soit $0 \longrightarrow X' \longrightarrow X \longrightarrow X'' \longrightarrow 0$ une suite exacte de systemes projectifs $J$-adiques noetheriens, munis de filtrations $J$-bonnes. Alors la suite des gradues associes est exacte :}
\[
0 \longrightarrow \mathrm{gr}(X') \longrightarrow \mathrm{gr}(X) \longrightarrow \mathrm{gr}(X'') \longrightarrow 0\;.
\]

\textbf{Preuve} : Comme les filtrations sont $J$-bonnes et les objets noetheriens, le lemme d'Artin-Rees 4.2.1 implique que les filtrations induites sont $J$-adiques. L'exactitude de la suite des gradues resulte alors de l'exactitude du foncteur gradue sur les suites strictes.

\medskip
\noindent\textbf{4.2.6.} \textbf{PROPOSITION.} \underline{Soit $X$ un systeme projectif $J$-adique noetherien. Alors il existe un entier $r \geq 0$ tel que, pour tout sous-systeme projectif strict $Y$ de $X$, le gradue $\mathrm{gr}(Y[r])$ soit un $\mathrm{gr}_J(A)$-module gradue.}

\textbf{Preuve} : Comme $X$ est noetherien, l'ensemble des sous-systemes projectifs stricts de $X$ est noetherien. Il existe donc un entier $r$ tel que, pour tout sous-systeme projectif strict $Y$, le systeme projectif $Y[r]$ soit $J$-adapte. On conclut par 2.2.4.

\subsection{Le lemme de Nakayama}

\medskip
\noindent\textbf{4.3.1.} \textbf{(Lemme de Nakayama).} \underline{Soit $M$ un objet d'une $A$-categorie abelienne $\cC$. On suppose que $M$ est $J$-adiquement complet et separe. Alors, si $M/JM = 0$, on a $M = 0$.}

\textbf{Preuve} : Si $M/JM = 0$, alors $M = JM$. Comme $M$ est $J$-adiquement complet et separe, on a $M = \varprojlim_n M/J^n M$. Or $M = JM$ implique $M/J^n M = 0$ pour tout $n$, donc $M = 0$.

\medskip
\noindent\textbf{4.3.2.} \textbf{COROLLAIRE.} \underline{Soit $u : M \longrightarrow N$ un morphisme entre objets $J$-adiquement complets et separes de $\cC$. Alors $u$ est un epimorphisme si et seulement si le morphisme induit $M/JM \longrightarrow N/JN$ est un epimorphisme.}

\textbf{Preuve} : Si $M/JM \longrightarrow N/JN$ est un epimorphisme, alors $\Coker(u)/J \cdot \Coker(u) = 0$. Comme $\Coker(u)$ est $J$-adiquement complet et separe (par exactitude du foncteur de completion sur les objets noetheriens), on a $\Coker(u) = 0$ d'apres 4.3.1.

\medskip
\noindent\textbf{4.3.3.} \textbf{COROLLAIRE.} \underline{Soit $M$ un objet $J$-adiquement complet et separe. Alors $M$ est de type fini (resp.\ noetherien) si et seulement si $M/JM$ est de type fini (resp.\ noetherien) en tant que $A/J$-module.}

\textbf{Preuve} : Si $M/JM$ est de type fini, il existe un morphisme $A^r \longrightarrow M$ tel que le morphisme induit $(A/J)^r \longrightarrow M/JM$ soit surjectif. D'apres 4.3.2, $A^r \longrightarrow M$ est surjectif, donc $M$ est de type fini. Le cas noetherien se demontre de meme.

\medskip
\noindent\textbf{4.3.4.} \textbf{PROPOSITION.} \underline{Soit $M$ un objet $J$-adiquement complet et separe de $\cC$. Alors $M$ est noetherien si et seulement si $M/JM$ est noetherien et les noyaux des morphismes $J^n \otimes_A M \longrightarrow J^n M$ sont de type fini.}

\textbf{Preuve} : Si $M$ est noetherien, alors $M/JM$ est noetherien, et les morphismes $J^n \otimes_A M \longrightarrow J^n M$ sont surjectifs (d'apres 4.2.2). Reciproquement, si $M/JM$ est noetherien, il existe un morphisme surjectif $A^r \longrightarrow M$. Il suffit alors de montrer que le noyau est de type fini, ce qui resulte de l'hypothese sur les noyaux.

\section{Systemes projectifs J-adiques et AR-J-adiques noetheriens}

\medskip
\noindent\textbf{5.0.} Dans cette section, $\cC$ designe une $A$-categorie abelienne, et $J$ un ideal de $A$ verifiant la condition (AR) de 1.1.8 (ii). On suppose de plus que $\cC$ possede des sommes directes denombrables.

\subsection{Definitions et premieres proprietes}

\medskip
\noindent\textbf{5.1.1.} On dit qu'un systeme projectif $X = (X_n)$ est \underline{AR-$J$-adique noetherien} (resp.\ \underline{AR-$J$-adique artinien}) s'il est AR-isomorphe a un systeme projectif strict $J$-adique noetherien (resp.\ artinien). On notera $\mathrm{AR\text{-}J\text{-}adn}(\cC)$ (resp.\ $\mathrm{AR\text{-}J\text{-}adf}(\cC)$) la sous-categorie pleine de $\Hom(\bN^\circ, \cC)$ engendree par les images des systemes projectifs AR-$J$-adiques noetheriens (resp.\ artiniens).

\medskip
\noindent\textbf{5.1.2.} \textbf{LEMME.} \underline{Soit $X$ un systeme projectif AR-$J$-adique noetherien. Alors, pour tout entier $r \geq 0$, le systeme projectif $X[r]$ est AR-$J$-adique noetherien.}

\textbf{Preuve} : Il existe un systeme projectif strict $J$-adique noetherien $Y$ et un AR-isomorphisme $Y \longrightarrow X$. Alors $Y[r] \longrightarrow X[r]$ est un AR-isomorphisme, et $Y[r]$ est $J$-adique noetherien d'apres 3.1.5.

\medskip
\noindent\textbf{5.1.3.} \textbf{DEFINITION.} \underline{Soit $X$ un systeme projectif indexe par $\bN$ a valeurs dans $\cC$. On dit qu'il est AR-$J$-adique noetherien (resp.\ artinien) s'il est AR-$J$-adique, et si ses composantes sont des objets noetheriens (resp.\ artiniens).}

On notera $\mathrm{AR\text{-}J\text{-}adn}(\cC)$ (resp.\ $\mathrm{AR\text{-}J\text{-}adf}(\cC)$) la sous-categorie pleine de $\underline{\Hom}_{AR}(\bN^\circ, \cC)$ (2.4) engendree par les images des systemes projectifs AR-$J$-adiques noetheriens (resp.\ artiniens).

\medskip
\noindent\textbf{5.1.4.} Il est alors clair que les fonctions canoniques
\begin{align*}
p_{AR}^n : J\text{-}\mathrm{adn}(\cC) &\longrightarrow \mathrm{AR\text{-}J\text{-}adn}(\cC) \\
p_{AR}^f : J\text{-}\mathrm{adf}(\cC) &\longrightarrow \mathrm{AR\text{-}J\text{-}adf}(\cC)
\end{align*}
obtenus par restriction de $p_{AR}$ (2.4) sont des equivalences de categories.

Rappelons brievement l'enonce et la demonstration du theoreme de Hilbert :

\medskip
\noindent\textbf{THEOREME 5.1.4 (Theoreme de Hilbert).} \underline{Soient $A$ un anneau commutatif, $\cC$ une $A$-categorie abelienne, $M$ un objet de type fini (resp.\ noetherien) de $\cC$.}

\underline{Pour toute $A$-algebre de type fini $B$ graduee par $\bN$, le $(B,\cC)$-module gradue $B \otimes_A M$ est de type fini (resp.\ noetherien).}

(N.B. Si $B = \bigoplus_{n \geq 0} B_n$, $B \otimes_A M$ designe l'objet gradue $(B_n \otimes_A M)_{n \in \bN}$.)

\textbf{Preuve} : On se ramene immediatement au cas ou $B$ est l'anneau de polynomes $A[T]$, muni de sa graduation habituelle. Le $(B,\cC)$-module $A[T] \otimes_A M = (M_n)_{n \in \bN}$ s'obtient alors en posant $M_n = M$ pour tout $n$, $T$ envoyant, pour tout entier $n$, $M_n$ dans $M_{n+1}$ par l'identite de $M$.

Soit $N$ un sous-$(A[T],\cC)$-module gradue de $A[T] \otimes_A M$. Les composantes $(N_n)_{n \in \bN}$ de $N$ definissent une suite croissante :
\[
N_0 \subset N_1 \subset \cdots \subset N_n \subset \cdots
\]
de sous-objets de $M$. Supposons maintenant que $M$ soit noetherien (resp.\ de type fini) et soit
\[
N_1 \subset N_2 \subset \cdots \subset N_p \subset \cdots
\]
une suite croissante de sous-$(A[T],\cC)$-modules gradues de $A[T] \otimes_A M$, qui dans la premiere hypothese (lorsque $M$ est de type fini) a pour reunion $A[T] \otimes_A M$.

Considerons alors le tableau a double entree ci-dessous :
\[
\begin{array}{ccccccccc}
N_0^1 & \subset & N_0^2 & \subset & \cdots\cdots\cdots\cdots & \subset & N_0^p & \subset & \cdots\cdots\cdots \\[-2pt]
\cap & & \cap & & & & \cap \\[-2pt]
N_1^1 & & N_1^2 & & & & N_1^p \\[-2pt]
\cap & & \cap & & & & \cap \\[-2pt]
\vdots & & \vdots & & & & \vdots \\[-2pt]
\cap & & \cap & & & & \cap \\[-2pt]
N_n^1 & \subset & N_n^2 & \subset & \cdots\cdots\cdots\cdots & \subset & N_n^p & \cdots\cdots\cdots \\[-2pt]
\cap & & \cap & & & & \cap \\[-2pt]
\vdots & & \vdots & & & & \vdots
\end{array}
\]

Il est clair par hypothese sur $M$ qu'il existe un couple $(p_0, n_0)$ tel que $N_{n_0}^{p_0}$ soit maximal (resp. egal a $M$).

Si $p_1$ est un entier $\geq p_0$ tel que les $N_n^{p_1}$ ($n \leq n_0$) soient les plus grands elements des suites $(N_n^p)_{p \in \bN}$, il est clair que $N_n^p = N_n^{p_1}$ des que $p \geq p_1$, ce qui permet de conclure dans les deux cas.

\medskip
\noindent\textbf{COROLLAIRE 5.1.5.} \underline{Supposons de plus que la categorie $\cC$ possede des sommes directes denombrables. Alors pour tout objet $M$ de type fini (resp.\ noetherien) de $\cC$ et toute $A$-algebre de type fini $B$, le $(B,\cC)$-module $B \otimes_A M$ est de type fini (resp.\ noetherien).}

\textbf{Preuve} : On se ramene au cas ou $B = A[T]$. Soit $N$ un sous-$(A[T],\cC)$-module de $A[T] \otimes_A M$ (considere comme objet de $\cC$). Nous noterons $\mathrm{gr}(N)$ le gradue de $N$ pour la filtration induite par la filtration naturelle d'objet gradue de $A[T] \otimes_A M$ : $\mathrm{gr}(N)$ est un sous-$(A[T],\cC)$-module gradue du $(A[T],\cC)$-module gradue $A[T] \otimes_A M$.

Soit maintenant
\[
N^0 \subset N^1 \subset \cdots \subset N^p \subset \cdots
\]
une suite croissante de sous-$(A[T],\cC)$-modules de $A[T] \otimes_A M$, de reunion $A[T] \otimes_A M$ dans la premiere hypothese. Dans la premiere hypothese ($M$ de type fini), il existe un entier $i$ tel que $\mathrm{gr}(N^i) = A[T] \otimes_A M$, d'ou $N^i = A[T] \otimes_A M$. Dans la deuxieme hypothese ($M$ noetherien), la suite des $\mathrm{gr}(N^i)$ est stationnaire, donc la suite des $N^i$ egalement.

\medskip
\noindent\textbf{PROPOSITION 5.1.6.} \underline{Soient $A$ un anneau commutatif noetherien, $J$ un ideal de $A$, $\cC$ une $A$-categorie abelienne. Pour tout systeme projectif $J$-adique noetherien} (5.1.1) $X$, \underline{le gradue $\mathrm{gr}_S(X)$ est un $(\mathrm{gr}_J(A),\cC)$-module noetherien.}

\textbf{Preuve} : Comme $\mathrm{gr}_J(A)$ est une $A$-algebre de type fini, il resulte du theoreme de Hilbert (5.1.4) que $\mathrm{gr}_J(A) \otimes_A X_0$ est un $(\mathrm{gr}_J(A),\cC)$-module gradue noetherien. L'assertion resulte alors de l'existence de l'epimorphisme (4.2.4) :
\[
\mathrm{gr}_J(A) \otimes_A X_0 \longrightarrow \mathrm{gr}(X)\;.
\]

\subsection{Structure de $J\text{-}\mathrm{adn}(\cC)$ et $\mathrm{AR\text{-}J\text{-}adn}(\cC)$}

\medskip
\noindent\textbf{5.2.} Dans cette section, nous designerons par $\cE_\cC$, ou $\cE$ s'il n'y a pas de confusion possible, la sous-categorie pleine de $\underline{\Hom}(\bN^\circ, \cC)$ dont les objets sont les systemes projectifs AR-$J$-adiques noetheriens.

\medskip
\noindent\textbf{PROPOSITION 5.2.1.} \underline{La categorie $\cE_\cC$ est stable par noyaux et conoyaux dans $\underline{\Hom}(\bN^\circ, \cC)$. Autrement dit, $\cE_\cC$ est une categorie abelienne et le foncteur d'inclusion :}
\[
\cE_\cC \longrightarrow \underline{\Hom}(\bN^\circ, \cC)
\]
\underline{est exact.}

\textbf{Preuve} : Soit $u : X \longrightarrow Y$ un morphisme de $\cE_\cC$. Nous allons voir successivement que le conoyau, l'image et le noyau de $u$ sont AR-$J$-adiques noetheriens. (Bien entendu la demonstration du fait que l'image de $u$ est AR-$J$-adique noetherienne est uniquement un intermediaire technique.)

a) \underline{$\Coker(u)$ est AR-$J$-adique noetherien.} Comme $X$ verifie (MLAR) et $Y$ est AR-$J$-adique, $\Coker(u)$ est AR-$J$-adique (3.2.4 (i)).

b) \underline{$\Im(u)$ est AR-$J$-adique noetherien.}

Grace a la suite exacte :
\[
0 \longrightarrow \Im(u) \longrightarrow Y \longrightarrow \Coker(u) \longrightarrow 0\;,
\]
l'assertion resultera du lemme 5.2.1.1 ci-dessous.

\medskip
\noindent\textbf{LEMME 5.2.1.1.} \underline{Soit}
\[
0 \longrightarrow L \longrightarrow M \longrightarrow N \longrightarrow 0
\]
\underline{une suite exacte dans la categorie $\underline{\Hom}(\bN^\circ, \cC)$. Si $M$ et $N$ sont AR-$J$-adiques noetheriens, $L$ est AR-$J$-adique noetherien.}

Il suffit evidemment de voir qu'il est AR-$J$-adique. On ne ramene de la meme maniere que dans la demonstration de (3.2.4) au cas ou $M$ et $N$ sont stricts, puis meme $J$-adiques noetheriens. Dans ce cas, il resulte de 3.1.3 (ii) que $L$ est strict. Le systeme projectif $L$ est strict, adapte a $J$ et son gradue, comme sous-objet gradue de $\mathrm{gr}(M)$, est noetherien : $\mathrm{gr}(M)$, quotient de $\mathrm{gr}_J(A) \otimes_A M_0$, est noetherien (5.1.6). L'assertion resulte alors du lemme d'Artin-Rees (4.2.6).

c) \underline{$\Ker(u)$ est AR-$J$-adique noetherien.} Il suffit d'appliquer (5.2.1.1) a la suite exacte evidente ci-dessous :
\[
0 \longrightarrow \Ker(u) \longrightarrow X \longrightarrow \Im(u) \longrightarrow 0\;.
\]

\medskip
\noindent\textbf{PROPOSITION 5.2.2.} \underline{Les objets de la categorie abelienne $\mathrm{AR\text{-}J\text{-}adn}(\cC)$ sont noetheriens.}

\textbf{Preuve} : On se ramene immediatement a voir que si l'on a un objet $J$-adique noetherien $X$, une suite d'objets $J$-adiques noetheriens $(Y^n)_{n \in \bN}$ et des morphismes de $\underline{\Hom}(\bN^\circ, \cC)$
\[
v_n : Y^n \longrightarrow X
\]
tels que
\begin{itemize}
\item[(i)] les $v_n$ deviennent des monomorphismes dans $\mathrm{AR\text{-}J\text{-}adn}(\cC)$,
\item[(ii)] pour tout entier $n$, $v_n$ se factorise a travers $v_{n+1}$ dans $\mathrm{AR\text{-}J\text{-}adn}(\cC)$,
\end{itemize}
alors les $v_n$ restent stationnaires, a isomorphisme pres, pour $n$ assez grand, dans $\mathrm{AR\text{-}J\text{-}adn}(\cC)$.

Or, dire que $v_n$ se factorise a travers $v_{n+1}$ dans $\mathrm{AR\text{-}J\text{-}adn}(\cC)$, signifie qu'il existe un entier $r$ et un morphisme
\[
f : Y^n[r] \longrightarrow Y^{n+1}
\]
tels que, $a$ designant le morphisme canonique
\[
Y^n[r] \longrightarrow Y^n\;,
\]
l'egalite ci-dessous ait lieu :
\[
v_n \circ a = v_{n+1} \circ f\;.
\]

Comme le morphisme $a$ est strict, on en deduit en particulier que l'image de $v_n$ est contenue dans celle de $v_{n+1}$. Quitte a remplacer chaque $Y^n$ par $\Im(v_n)$, on est ramene a voir que toute suite croissante de sous-systemes projectifs stricts de $X$ est stationnaire. Or, si $(Z^n)_{n \in \bN}$ est une telle suite, vu que $\mathrm{gr}(X)$ est noetherien, la suite croissante des $\mathrm{gr}(Z^n)$ est stationnaire, donc aussi celle des $Z^n$.

\medskip
\noindent\textbf{Resume.} Les resultats precedents se resument comme suit :
\begin{itemize}
\item[(i)] La categorie $\cE_\cC$ des systemes projectifs AR-$J$-adiques noetheriens est une sous-categorie abelienne epaisse de $\underline{\Hom}(\bN^\circ, \cC)$.
\item[(ii)] Le foncteur canonique $Q : \underline{\Hom}(\bN^\circ, \cC) \longrightarrow \underline{\Hom}(\bN^\circ, \cC) / \underline{\Hom}(\bN^\circ, \cC)_{\mathrm{AR\text{-}nul}}$ induit une equivalence de categories
\[
\cE_\cC / \cE_\cC \cap \underline{\Hom}(\bN^\circ, \cC)_{\mathrm{AR\text{-}nul}} \xrightarrow{\sim} \mathrm{AR\text{-}J\text{-}adn}(\cC)\;.
\]
\item[(iii)] Les objets de $\mathrm{AR\text{-}J\text{-}adn}(\cC)$ sont noetheriens.
\end{itemize}

Ces proprietes permettent d'etendre a la categorie $\mathrm{AR\text{-}J\text{-}adn}(\cC)$ les resultats classiques sur les categories abeliennes de modules noetheriens, notamment en ce qui concerne les foncteurs derives et la dualite.

\bigskip
\hrule
\bigskip

\noindent\textbf{Fin de l'Expose V.}


\clearpage
% ===== Source block: SGA 5 pages 251 300 =====
\phantomsection
\addcontentsline{toc}{section}{SGA 5 pages 251 300}
\section*{SGA 5 pages 251 300}
\setcounter{page}{251}

\noindent
\textbf{THEOR\`EME 5.2.3.} \underline{Les cat\'egories} $\mathrm{J\text{-}adn}(\cC)$ \underline{et} $\mathrm{AR\text{-}J\text{-}adn}(\cC)$ \underline{sont ab\'eliennes} \underline{et noeth\'eriennes.}

En ce qui concerne $\mathrm{AR\text{-}J\text{-}adn}(\cC)$, c'est ce que nous venons de voir. Quant \`a $\mathrm{J\text{-}adn}(\cC)$, elle lui est \'equivalente.

\bigskip

\noindent
\textbf{PROPOSITION 5.2.4.} (\underline{stabilit\'es par extension}). \underline{Soit}
\[
0 \longrightarrow X \longrightarrow Y \longrightarrow Z \longrightarrow 0
\]
\underline{une suite exacte de} $\underline{\Hom}(\bN^{\circ}, \cC)$, \underline{avec} $J^{n+1}Y_n = 0$ \underline{pour tout} $n$. \underline{On suppose que :}
\begin{itemize}
\item[(i)] $Z$ \underline{est} $\mathrm{AR\text{-}J\text{-}adique}$,
\item[(ii)] $X$ \underline{est} $\mathrm{AR\text{-}J\text{-}adique}$ artinien.
\end{itemize}
\underline{Alors} $Y$ \underline{est} $\mathrm{AR\text{-}J\text{-}adique}$.

\smallskip
\noindent
\underline{Preuve} : On se ram\`ene comme dans la d\'emonstration de (3.2.4) au cas o\`u $Y$ et $Z$ sont stricts et $Z$ est J-adique. Alors (3.1.3 (ii)) $X$ est strict. Il existe alors un entier $r$ tel que le syst\`eme projectif $(X_{n+r}/J^{n+1}X_{n+r})_{n \in \bN}$ soit J-adique artinien. Comme $X$ en est un quotient et lui est AR-isomorphe, on est ramen\'e \`a prouver le lemme suivant :

\bigskip

\noindent
\textbf{LEMME 5.2.4.1.} \underline{Soit}
\[
0 \longrightarrow T \longrightarrow X \longrightarrow Y \longrightarrow Z \longrightarrow 0
\]
\underline{une suite exacte de} $\underline{\Hom}(\bN^{\circ}, \cC)$, \underline{avec} $J^{n+1}Y_n = 0$ \underline{pour tout} $n$. \underline{Si} $X$ \underline{est} J-adique artinien, $Z$ \underline{est} J-adique et $T$ \underline{est} AR-nul, \underline{alors} $Y$ \underline{est} AR-J-adique.

Pour montrer le lemme, on va d'abord voir que, pour tout entier $n$, le syst\`eme projectif \'evident $(Y_{n+k}/J^{n+1}Y_{n+k})_{k \in \bN}$ est essentiellement constant. D\'esignant par $S_n^k$ le noyau du morphisme \'evident
\[
X_{n+k}/J^{n+1}X_{n+k} \longrightarrow Y_{n+k}/J^{n+1}Y_{n+k},
\]

\newpage

consid\'erons le diagramme commutatif exact \'evident ci-dessous :
\[
\begin{array}{ccccccccc}
0 & \longrightarrow & S_n^k & \longrightarrow & X_{n+k}/J^{n+1}X_{n+k} & \longrightarrow & Y_{n+k}/J^{n+1}Y_{n+k} & \longrightarrow & Z_{n+k}/J^{n+1}Z_{n+k} \longrightarrow 0 \\
  &                 & \downarrow &              & \downarrow{\wr} &              & \downarrow &              & \downarrow{\wr} \\
0 & \longrightarrow & T_n & \longrightarrow & X_n & \longrightarrow & Y_n & \longrightarrow & Z_n \longrightarrow 0\,. \\
\end{array}
\]

L'entier $n$ \'etant fix\'e, les objets $S_n^k$ s'identifient \`a des sous-objets de $T_n$ et forment une suite d\'ecroissante. Comme $T_n$ est artinien, il existe donc un entier $k_0$ tel que $S_n^k = S_n^{k_0}$ pour tout entier $k \geq k_0$, d'o\`u il r\'esulte aussit\^ot que les morphismes canoniques
\[
Y_{n+k}/J^{n+1}Y_{n+k} \longrightarrow Y_{n+k_0}/J^{n+1}Y_{n+k_0}
\]
soient des isomorphismes.

Ceci dit, d\'esignons par $Y'_n$ le syst\`eme projectif d\'efini par
\[
Y'_n = \varprojlim_{r} Y_{n+r}/J^{n+1}Y_{n+r}
\]
pour les objets, et de fa\c con \'evidente pour les morphismes. Le syst\`eme projectif $Y'$ est J-adique. De plus on obtient ``par passage \`a la limite'' un diagramme commutatif exact du type ci-dessous :
\[
\begin{array}{ccccccccc}
0 & \longrightarrow & S & \longrightarrow & X & \longrightarrow & Y' & \longrightarrow & Z \longrightarrow 0 \\
  &                 & \downarrow &              & \downarrow{\mathrm{id}} &     & \downarrow &              & \downarrow{\mathrm{id}} \\
0 & \longrightarrow & T & \longrightarrow & X & \longrightarrow & Y & \longrightarrow & Z \longrightarrow 0\,. \\
\end{array}
\]

Comme $T$ est AR-nul, le lemme du serpent appliqu\'e \`a ce diagramme montre que le morphisme \'evident $Y' \rightarrow Y$ est un AR-isomorphisme, d'o\`u le fait que $Y$ est AR-J-adique.

\newpage

\noindent
\textbf{COROLLAIRE 5.2.5.} \underline{Soit}
\[
0 \longrightarrow X \longrightarrow Y \longrightarrow Z \longrightarrow 0
\]
\underline{une suite exacte de} $\underline{\Hom}(\bN^{\circ}, \cC)$, \underline{avec} $J^{n+1}Y_n = 0$ \underline{pour tout} $n$. \underline{On suppose que :}
\begin{itemize}
\item[(i)] $X$ \underline{est} $\mathrm{AR\text{-}J\text{-}adique}$ artinien (resp. artinien et noeth\'erien),
\item[(ii)] $Z$ \underline{est} $\mathrm{AR\text{-}J\text{-}adique}$ artinien (resp. noeth\'erien).
\end{itemize}
\underline{Alors} $Y$ \underline{est} $\mathrm{AR\text{-}J\text{-}adique}$ artinien (resp. noeth\'erien).

\bigskip

\noindent
\textbf{5.3.} \underline{Syst\`emes projectifs AR-J-adiques noeth\'eriens et $\partial$-foncteurs.}

Soient $\cC$ et $\cD$ deux A-cat\'egories ab\'eliennes et $F : \cC \rightarrow \cD$ un foncteur additif. Le foncteur $F$ se prolonge de fa\c con claire en un foncteur, not\'e encore $F$,
\[
F : \underline{\Hom}(\bN^{\circ}, \cC) \longrightarrow \underline{\Hom}(\bN^{\circ}, \cD)\,,
\]
qui est additif et ``commute aux translations''.

\bigskip

\noindent
\textbf{PROPOSITION 5.3.1.} \underline{Soient} $\cC$ \underline{et} $\cD$ \underline{deux} A-\underline{cat\'egories ab\'eliennes et}
\[
T = (T^i)_{i \in \bN}
\]
\underline{un} $\partial$-\underline{foncteur exact} A-\underline{lin\'eaire de} $\cC$ \underline{dans} $\cD$.

\underline{On suppose que :}
\begin{itemize}
\item[(i)] \underline{Pour tout objet noeth\'erien} $M$ \underline{de} $\cC$ \underline{annul\'e par une puissance de} $J$ \underline{et tout entier} $n$, $T^n(M)$ \underline{est noeth\'erien.}
\item[(ii)] \underline{Pour toute} A-\underline{alg\`ebre gradu\'ee de type fini} $B$, \underline{annul\'ee par} $J$, \underline{et tout} $(B,\cC)$-\underline{module gradu\'e noeth\'erien} $N$, \underline{les} $(B,\cD)$-\underline{modules gradu\'es} $T^n(N) (n \in \bN)$ \underline{sont noeth\'eriens.}
\end{itemize}

\hspace*{\parindent}
\underline{Alors pour tout objet} $\mathrm{AR\text{-}J\text{-}adique}$ noeth\'erien $X$ \underline{de} $\underline{\Hom}(\bN^{\circ}, \cC)$, \underline{les} $T^n(X)$ \underline{sont} $\mathrm{AR\text{-}J\text{-}adiques}$ noeth\'eriens.

\newpage

\noindent
\underline{Preuve} : Comme les foncteurs $T^n$ commutent aux translations, il suffit de montrer l'assertion lorsque $X$ est J-adique noeth\'erien. Dans ce cas, comme $\mathrm{gr}_J(X)$ est un $\mathrm{gr}_J(A)$-module gradu\'e noeth\'erien (5.1.6), il r\'esulte de l'hypoth\`ese (ii) et d'une variante d'un corollaire du th\'eor\`eme de SHIH (cf. Appendice A.3) que pour tout entier $n$ les assertions suivantes sont v\'erifi\'ees :
\begin{itemize}
\item[(a)] $T^n(X)$ satisfait (MLAR),
\item[(b)] $\mathrm{gr}_J\,T^n(X)$ est un $(\mathrm{gr}_J(A), \cD)$-module gradu\'e noeth\'erien.
\end{itemize}

Le ``lemme d'Artin-Rees'' (4.2.6) permet de conclure.

\bigskip

\noindent
\textbf{COROLLAIRE 5.3.2.} \underline{Le} $\partial$-\underline{foncteur}
\[
T = (T^i)_{i \in \bN} : \underline{\Hom}(\bN^{\circ}, \cC) \longrightarrow \underline{\Hom}(\bN^{\circ}, \cD)
\]
\underline{induit un} $\partial$-\underline{foncteur exact} A-\underline{lin\'eaire}
\[
\cE_{\cC} \longrightarrow \cE_{\cD}\,,
\]
\underline{commutant aux translations, donc par passage au quotient un nouveau} $\partial$-\underline{foncteur exact} A-\underline{lin\'eaire, not\'e encore} $(T^i)_{i \in \bN}$,
\[
\mathrm{AR\text{-}J\text{-}adn}(\cC) \longrightarrow \mathrm{AR\text{-}J\text{-}adn}(\cD)\,.\hspace*{3cm}
\]

\newpage

\begin{center}
\underline{Appendice : le th\'eor\`eme de SHIH.}
\end{center}

Le but du pr\'esent appendice est de donner une d\'emonstration du th\'eor\`eme de SHIH sous la forme utilis\'ee dans la preuve de (5.3.1). Cette forme \'etant plus g\'en\'erale que celle figurant dans $\mathrm{EGA}\,\cO_{\mathrm{III}}$, une d\'emonstration directe, adapt\'ee de celle de loc. cit., s'est r\'ev\'el\'ee n\'ecessaire.

\smallskip

\underline{Terminologie} : Soit $r_0$ un entier $\geq 1$. La donn\'ee, pour tout entier $r \geq r_0$, d'un objet bigradu\'e $E_r = (E_r^{pq})_{p \in \bZ, q \in \bZ}$, muni d'une diff\'erentielle $d_r$ de bidegr\'e $(r, 1-r)$, et pour tout $r$ d'un isomorphisme
\[
H(E_r) \xrightarrow{\;\sim\;} E_{r+1}\,,
\]
sera appel\'ee dans la suite \underline{filtre spectral}. On d\'efinit de fa\c con \'evidente la notion de morphisme de filtres spectraux.

\bigskip

\noindent
\textbf{A.1.} \underline{Rappels sur les suites spectrales associ\'ees aux} $\partial$-\underline{foncteurs.}

Soit $\cC$ une cat\'egorie ab\'elienne et $X$ un objet de $\cC$, muni d'une filtration d\'ecroissante $(F^p X)_{p \in \bZ}$. Pour tout $\partial$-foncteur exact de $\cC$ dans une cat\'egorie ab\'elienne $\cD$, la m\'ethode dite ``des couples exacts'' (cf. par ex. Maclane, Homology, p. 336) permet de d\'ecrire une suite spectrale :
\[
\mathbb{E}(X) : E_1^{pq} = T^{p+q}(F^p X / F^{p+1} X) \Rightarrow E^n = T^n(X)\,,
\]
pour laquelle la filtration de $E^n$ est d\'efinie comme suit :
\[
F^p(E^n) = \Im(T^n(F^p X) \longrightarrow T^n(X)) = \Ker(T^n(X) \longrightarrow T^n(X/F^p X))\,.\hspace*{1cm}
\]

De plus, si la filtration de $X$ est finie, la suite spectrale ainsi obtenue est faiblement convergente ($\mathrm{EGA}\,\cO_{\mathrm{III}}$ 11.1.3).

\newpage

\noindent
\underline{Fonctorialit\'e.} Si $f : X \rightarrow Y$ est un morphisme d'objets filtr\'es de $\cC$, on construit de fa\c con fonctorielle un morphisme de suites spectrales :
\[
\mathbb{E}(f) : \mathbb{E}(X) \longrightarrow \mathbb{E}(Y)\,,
\]
coincidant avec $T^{p+q}(\mathrm{gr}^p f)$ (resp. $T^n(f)$) sur $E_1^{pq}(X)$ (resp. $E^n(X)$).

En particulier, si $u : X \rightarrow X$ est un endomorphisme de $X$ tel que $u(F^p X) \subset F^{p+h}(X)$, pour un entier $h$ ind\'ependant de $p$, $\mathbb{E}(u)$ est un endomorphisme de bidegr\'e $(h, -h)$ du filtre spectral associ\'e \`a $\mathbb{E}(X)$, commutant avec les diff\'erentielles.

\bigskip

\noindent
\textbf{A.2.} \underline{Cas des syst\`emes projectifs.} Soit maintenant $X = (X_n)_{n \in \bZ}$ un syst\`eme projectif d'objets d'une cat\'egorie ab\'elienne $\cC$, tel que
\[
X_i = 0 \qquad (i < 0)\,.\hspace*{5cm}
\]

Si $T$ d\'esigne un $\partial$-foncteur exact de $\cC$ dans une cat\'egorie ab\'elienne $\cD$, $T$ se prolonge de fa\c con naturelle en un $\partial$-foncteur exact, not\'e encore $T$ :
\[
\underline{\Hom}(\bZ^{\circ}, \cC) \longrightarrow \underline{\Hom}(\bZ^{\circ}, \cD)\,.\hspace*{3cm}
\]

D'o\`u, munissant $X$ de sa filtration naturelle de syst\`eme projectif ($\mathrm{EGA}\,\cO_{\mathrm{III}}$ 13.4.1), une suite spectrale :
\[
\mathbb{E}(X) = \mathbb{E}(X_n)_{n \in \bZ}\,,\hspace*{5cm}
\]
o\`u les $\mathbb{E}(X_n)$ d\'esignent les suites spectrales associ\'ees \`a $T$ et aux $X_n$ munis de la filtration induite.

\smallskip

\noindent
\underline{Propri\'et\'es de} $\mathbb{E}(X)$.

a) Pour tout entier $n$, la filtration de l'aboutissement $E^n(X)$ dans la suite spectrale $\mathbb{E}(X)$ est identique \`a sa filtration naturelle de syst\`eme projectif, lorsque $X$ est \underline{strict}.

En effet, d\'esignant par la lettre $F$ (comme dans A.1) la premi\`ere

\newpage

filtration, on a :
\[
F^p(E^n(X)) = \Im(T^n(F^p X) \rightarrow T^n(X))\,.\hspace*{4cm}
\]
Or il r\'esulte de la d\'efinition de la filtration naturelle de syst\`eme projectif que $(F^p X)_k = X_k$ si $k \geq p$ et $= 0$ si $k < p$. Par suite :
\[
F^p(E^n(X))_k = \Im(T^n(X_k) \rightarrow T^n(X_k)) = T^n(X_k) \quad (k \geq p)
\]
et
\[
F^p(E^n(X))_k = \Im(0 \rightarrow T^n(X_k)) = 0 \quad (k < p)\,.
\]

b) Soit $S$ un anneau op\'erant sur le syst\`eme projectif $X$ de fa\c con compatible avec sa filtration naturelle de syst\`eme projectif (i.e. $S$ op\`ere sur les $X_n$ de fa\c con compatible avec les morphismes de transition et $F^p S = \bigoplus_{i \geq p} S^i$ si $S$ est gradu\'e). Alors, pour tout entier $r \geq 1$, les objets bigradu\'es
\[
E_r(X) = (E_r^{pq}(X))\hspace*{5cm}
\]
sont canoniquement munis d'une structure de $(\mathrm{gr}^{\bullet}(S), \cD)$-module gradu\'e. Il en est de m\^eme par passage \`a la limite pour :
\[
\cE_r^{(n)}(X) = (\cE_r^{p,q}(X))_{p+q=n}\,.\hspace*{5cm}
\]

L'anneau $S$, op\'erant sur $X$ de fa\c con compatible avec sa filtration de syst\`eme projectif, op\`ere de m\^eme sur $T^n(X)$ pour tout entier $n$. Par suite le gradu\'e strict $\mathrm{gr}_S(X)$ est canoniquement muni d'une structure de $(\mathrm{gr}^{\bullet}(S), \cC)$-module gradu\'e et de m\^eme, lorsque $T^n(X)$ v\'erifie (ML), le gradu\'e strict $\mathrm{gr}_S^{\bullet}\,T^n(X)$ est canoniquement muni d'une structure de $(\mathrm{gr}^{\bullet}(S), \cD)$-module gradu\'e.

\bigskip

\noindent
\textbf{THEOR\`EME A.3.} \underline{Supposons que pour un entier} $n$ \underline{les} $(\mathrm{gr}^{\bullet}(S), \cD)$-\underline{modules gradu\'es}
\[
\cE_1^{(n)}(X) = T^n(\mathrm{gr}_S^{\bullet}(X)) \quad \text{et} \quad \cE_1^{(n-1)}(X) = T^{n-1}(\mathrm{gr}_S^{\bullet}(X))
\]
\underline{soient noeth\'eriens. Alors :}
\begin{itemize}
\item[(i)] $T^n(X)$ \underline{v\'erifie} (MLAR),
\item[(ii)] $\mathrm{gr}_S^{\bullet}\,T^n(X)$ \underline{est un} $(\mathrm{gr}^{\bullet}(S), \cD)$-\underline{module noeth\'erien.}
\end{itemize}

\smallskip

\noindent
\underline{Preuve} :

i) Pour tout couple $(m,r)$ d'entiers,
\[
\cE_r^{(m)}(X) = (\cE_r^{p,m-p}(X))_{p \in \bN}\hspace*{3cm}
\]
est un sous-$(\mathrm{gr}^{\bullet}(S), \cD)$-module gradu\'e de $\cE_1^{(m)}(X)$. En particulier, pour tout entier $m$, les $\cE_r^{(m)}(X)$ forment une suite croissante index\'ee par $r$ de sous-$(\mathrm{gr}^{\bullet}(S), \cD)$-modules gradu\'es de $\cE_1^{(m)}(X)$. Utilisant le fait que $\cE_1^{(n)}(X)$ et $\cE_1^{(n-1)}(X)$ sont des $(\mathrm{gr}^{\bullet}(S), \cD)$-modules gradu\'es noeth\'eriens, on en d\'eduit que les suites

\newpage

$(\cE_r^{(n)}(X))_{r \geq 1}$ et $(\cE_r^{(n-1)}(X))_{r \geq 1}$ sont stationnaires. Il existe donc un entier $r_0$ tel que, pour tout entier $r \geq r_0$, on ait :
\[
\cE_r^{(n)}(X) = \cE_{r_0}^{(n)}(X) \quad \text{et} \quad \cE_r^{(n-1)}(X) = \cE_{r_0}^{(n-1)}(X)\,.\hspace*{1cm}
\]

Utilisant une caract\'erisation \'equivalente de la condition de Mittag-Leffler, on voit que la suite spectrale $\mathbb{E}(X)$ est faiblement convergente dans les degr\'es $n$ et $n-1$ (EGA $\cO_{\mathrm{III}}$ 11.2.1). Il en r\'esulte que la filtration de $E^n(X) = T^n(X)$ est s\'epar\'ee et que $\cE_{\infty}^{(n)}(X)$ est un sous-objet de $\cE_{r_0}^{(n)}(X)$. Comme ce dernier est noeth\'erien, $\cE_{\infty}^{(n)}(X)$ est noeth\'erien.

Pour montrer que $T^n(X)$ v\'erifie (ML), il suffit de montrer que pour tout entier $p$, la suite des
\[
\Im(T^n(F^p X) \rightarrow T^n(X/F^{p+r} X))_{r \geq 1}
\]
est stationnaire. Or il r\'esulte de la faible convergence que, pour $r$ assez grand, cette image est le noyau d'un morphisme :
\[
T^n(X/F^p X) \longrightarrow E_{\infty}^{(n)}(X)\,,
\]
de sorte qu'elle est noeth\'erienne. La suite est donc bien stationnaire.

ii) R\'esulte de (i) et du fait que $\mathrm{gr}_S^{\bullet}\,T^n(X)$ est un sous-quotient de $\cE_{\infty}^{(n)}(X)$.

\bigskip

\noindent
\textbf{A.4.} \underline{Application aux} $\partial$-\underline{foncteurs.}

Soient $\cC$ et $\cD$ deux A-cat\'egories ab\'eliennes et $T = (T^i)_{i \in \bN}$ un $\partial$-foncteur exact A-lin\'eaire de $\cC$ dans $\cD$. Soient $X$ un objet de $\underline{\Hom}(\bN^{\circ}, \cC)$, $S$ une A-alg\`ebre gradu\'ee, $S$ op\'erant sur $X$ de fa\c con compatible avec les morphismes de transition. Supposons que :
\begin{itemize}
\item[(i)] $X$ est AR-J-adique noeth\'erien,
\item[(ii)] Pour tout objet noeth\'erien $M$ de $\cC$ annul\'e par une puissance de $J$ et tout entier $n$, $T^n(M)$ est noeth\'erien,
\item[(iii)] Pour toute A-alg\`ebre gradu\'ee de type fini $B$, annul\'ee par $J$, et tout $(B,\cC)$-module gradu\'e noeth\'erien $N$, les $(B,\cD)$-modules gradu\'es $T^n(N) (n \in \bN)$ sont noeth\'eriens.
\end{itemize}

\hspace*{\parindent}
\textbf{THEOR\`EME A.4.1.} \underline{Sous les hypoth\`eses pr\'ec\'edentes, les} $T^n(X)$ \underline{sont} $\mathrm{AR\text{-}J\text{-}adiques}$ \underline{noeth\'eriens.}

\smallskip

\noindent
\underline{Preuve} : On se ram\`ene par translation au cas o\`u $X$ est J-adique noeth\'erien. Alors $\mathrm{gr}_J^{\bullet}(X)$ est un $(\mathrm{gr}_J^{\bullet}(A), \cC)$-module gradu\'e noeth\'erien. Il r\'esulte alors de (iii) que les modules $\cE_1^{(n)}(X)$ et $\cE_1^{(n-1)}(X)$ sont noeth\'eriens. Le th\'eor\`eme A.3 permet de conclure.

\newpage

\begin{center}
\Large\textbf{EXPOSE VI}\\[0.5em]
\textbf{COHOMOLOGIE $\ell$-ADIQUE}
\end{center}

\bigskip

\noindent
\textbf{1.} \underline{Cat\'egories de faisceaux $\ell$-adiques constructibles.}

\noindent
\textbf{1.1.} \underline{G\'en\'eralit\'es.}

\noindent
\textbf{1.1.1.} Soient $X$ un sch\'ema et $\ell$ un nombre premier. On d\'esigne par
\[
\bZ_{\ell}\text{-}\mathrm{fc}(X)
\]
la cat\'egorie des $\bZ_{\ell}$-faisceaux constructibles sur $X$, d\'efinie dans (SGA 5 V). C'est une cat\'egorie ab\'elienne. On note
\[
\bQ_{\ell}\text{-}\mathrm{fc}(X)
\]
la cat\'egorie des $\bQ_{\ell}$-faisceaux constructibles, quotient de la pr\'ec\'edente par la sous-cat\'egorie \'epaisse des faisceaux de torsion. On a d\'efini \'egalement (loc. cit.) les cat\'egories
\[
\mathrm{AR}\text{-}\bZ_{\ell}\text{-}\mathrm{fc}(X) \quad \text{et} \quad \mathrm{AR}\text{-}\bQ_{\ell}\text{-}\mathrm{fc}(X)\,.\hspace*{1cm}
\]

\noindent
\textbf{1.1.2.} Rappelons que la donn\'ee d'un $\bZ_{\ell}$-faisceau constructible $\cF$ \'equivaut \`a celle d'un syst\`eme projectif $(\cF_n)_{n \in \bN}$ de faisceaux de $\bZ/\ell^{n+1}\bZ$-modules constructibles sur $X$, tel que pour tout $n$, le morphisme de transition $\cF_{n+1} \rightarrow \cF_n$ se factorise en un isomorphisme
\[
\cF_{n+1}/\ell^{n+1}\cF_{n+1} \xrightarrow{\;\sim\;} \cF_n\,.\hspace*{3cm}
\]

\noindent
\textbf{1.1.3.} Soient $\cF = (\cF_n)$ et $\cG = (\cG_n)$ deux $\bZ_{\ell}$-faisceaux constructibles. On a par d\'efinition :
\[
\Hom(\cF, \cG) = \varprojlim_{n} \Hom(\cF_n, \cG_n)\,.\hspace*{3cm}
\]

Pour tout entier $n$, le morphisme canonique
\[
\Hom(\cF, \cG) \longrightarrow \Hom(\cF_n, \cG_n)
\]
est surjectif et son noyau est l'ensemble des morphismes qui se factorisent par un objet de torsion.

\smallskip

\noindent
\textbf{1.1.4.} Soit $\cF = (\cF_n)_{n \in \bN}$ un faisceau $\ell$-adique constructible. Comme l'anneau $\bZ_{\ell}$ op\`ere de fa\c con \'evidente sur chaque faisceau $\cF_n$, il op\`ere \'egalement sur $\cF$. De plus, si $\cF$ et $\cG$ sont deux faisceaux $\ell$-adiques constructibles et
\[
u : \cF \longrightarrow \cG\hspace*{5cm}
\]
un morphisme, $u$ est compatible avec les op\'erations de $\bZ_{\ell}$ sur $\cF$ et $\cG$ respectivement. Ceci montre que $\ell\text{-}\mathrm{adc}(X)$ \underline{est canoniquement munie d'une structure de} $\bZ_{\ell}$-\underline{cat\'egorie ab\'elienne.}

Cela justifie la terminologie suivante :

\smallskip
\noindent
\underline{Terminologie} : On appellera dans la suite $\bZ_{\ell}$-faisceaux constructibles les faisceaux $\ell$-adiques constructibles.

On prendra garde de ne pas confondre cette notion avec celle de $(\bZ_{\ell})_X$-Module, qui ne sera pas utilis\'ee dans la suite.

La cat\'egorie $\ell\text{-}\mathrm{adc}(X)$ sera \'egalement not\'ee
\[
\bZ_{\ell}\text{-}\mathrm{fc}(X)\,.\hspace*{5cm}
\]

\newpage

\noindent
\textbf{1.2.} $\bZ_{\ell}$-\underline{faisceaux constants tordus constructibles.}

\noindent
\textbf{D\'efinition 1.2.1.} Un $\bZ_{\ell}$-faisceau constructible
\[
\cF = (\cF_n)_{n \in \bN}\hspace*{4cm}
\]
est dit \underline{constant tordu} si et seulement si les faisceaux $\cF_n$ sont localement constants sur $X_{\text{\'et}}$.

On appellera cat\'egorie des $\bZ_{\ell}$-faisceaux constants tordus constructibles la sous-cat\'egorie pleine de $\bZ_{\ell}\text{-}\mathrm{fc}(X)$ engendr\'ee par les $\bZ_{\ell}$-faisceaux constants tordus constructibles.

\bigskip

\noindent
\textbf{Exemple 1.2.2.} Le $\bZ_{\ell}$-faisceau $\bZ_{\ell}(1)$.

Supposons l'entier $\ell$ premier aux caract\'eristiques r\'esiduelles de $X$. Pour tout entier $n$, l'\'el\'evation \`a la puissance $\ell$ d\'efinit un morphisme de faisceaux
\[
\mathbb{G}_{\mathrm{m}} \longrightarrow \mathbb{G}_{\mathrm{m}}\,,
\]
d'o\`u par passage \`a la cohomologie un morphisme
\[
H^1(X, \mathbb{G}_{\mathrm{m}}) \longrightarrow H^1(X, \mathbb{G}_{\mathrm{m}})\,.\hspace*{2cm}
\]
On d\'esigne par $\bZ/\ell^{n+1}\bZ(1)$ le faisceau $\mu_{\ell^{n+1}}$ des racines $\ell^{n+1}$-\`emes de l'unit\'e sur $X$. Le syst\`eme projectif des $\mu_{\ell^{n+1}}$ d\'efinit un $\bZ_{\ell}$-faisceau constant tordu, not\'e $\bZ_{\ell}(1)$.

\smallskip

\noindent
\textbf{1.2.3.} Soit $\cF = (\cF_n)$ un $\bZ_{\ell}$-faisceau constant tordu constructible. Pour tout $x \in X$, la fibre g\'eom\'etrique $\cF_{\overline{x}}$ est un $\bZ_{\ell}$-module de type fini, muni d'une action continue du groupe fondamental $\pi_1(X, \overline{x})$. Le foncteur
\[
\cF \longmapsto \cF_{\overline{x}}\hspace*{5cm}
\]
est une \'equivalence de cat\'egories entre la cat\'egorie des $\bZ_{\ell}$-faisceaux constants tordus constructibles sur $X$ et la cat\'egorie des $\bZ_{\ell}$-modules de type fini munis d'une action continue de $\pi_1(X, \overline{x})$.

\smallskip

\noindent
\textbf{1.2.4.} Soit $M$ un $\bZ_{\ell}$-module de type fini. On d\'esigne par
\[
M_X\hspace*{5cm}
\]
le $\bZ_{\ell}$-faisceau constant tordu d\'efini par $M$. En particulier, on pose
\[
\bZ_{\ell}(r) = \bZ_{\ell}(1)^{\otimes r} \quad (r \in \bZ)\,.\hspace*{3cm}
\]

Pour $r \geq 0$, on a $\bZ_{\ell}(r) = \bZ_{\ell}(1)^{\otimes r}$, et pour $r < 0$, $\bZ_{\ell}(r) = \bZ_{\ell}(-r)^{\vee}$.

\newpage

\noindent
\textbf{1.2.4.1.} Soit $\cF = (\cF_n)$ un $\bZ_{\ell}$-faisceau constant tordu. Pour tout entier $i$, on d\'efinit
\[
H^i(X, \cF) = \varprojlim_{n} H^i(X, \cF_n)\,.\hspace*{3cm}
\]
Ces groupes sont des $\bZ_{\ell}$-modules de type fini si $X$ est un sch\'ema propre sur un corps s\'eparablement clos.

\bigskip

\noindent
\textbf{1.3.} \underline{G\'en\'eralit\'es sur les faisceaux $\ell$-adiques inversibles.}

\noindent
\textbf{1.3.1.} Soit $\cL$ un $\bZ_{\ell}$-faisceau constructible. Nous dirons qu'il est \underline{inversible} si pour tout entier $n$ son $n$\`eme composant $\cL_n$ est un $(\bZ/\ell^{n+1}\bZ)_X$-Module inversible, i.e. localement libre de rang un. Il revient au m\^eme de dire que $\cL$ est localement libre au sens du paragraphe pr\'ec\'edent et que les fibres g\'eom\'etriques des $\cL_n$ sont des groupes ab\'eliens cycliques non nuls.

Si $\cL$ est un $\bZ_{\ell}$-faisceau inversible, on convient de poser
\[
\cL^{-1} = \cH om(\cL, \bZ_{\ell})\,.\hspace*{4cm}
\]

Il est clair qu'on a un isomorphisme
\[
\cL \otimes \cL^{-1} \xrightarrow{\;\sim\;} \bZ_{\ell}\,,\hspace*{3cm}
\]
d\'eduit de l'isomorphisme analogue composant par composant (EGA $\cO_{\mathrm{I}}$ 5.4.3.1). Suivant alors (EGA $\cO_{\mathrm{I}}$ 5.4.4), on convient de poser $\cL^{\otimes 0} = \bZ_{\ell}$ et pour $n \geq 1$, $\cL^{\otimes (-n)} = (\cL^{-1})^{\otimes n}$. Avec ces notations, il existe alors un isomorphisme fonctoriel canonique
\[
\cL^{\otimes m} \otimes \cL^{\otimes n} \xrightarrow{\;\sim\;} \cL^{\otimes (m+n)}\hspace*{2cm}
\]
quels que soient les entiers rationnels $m$ et $n$. De plus cet isomorphisme jouit

\newpage

des propri\'et\'es d'associativit\'e habituelles.

En particulier, les objets gradu\'es
\[
(\cL^{\otimes i})_{i \in \bN} \quad \text{et} \quad (\cL^{\otimes i})_{i \in \bZ}\hspace*{2cm}
\]
sont des $\bZ_{\ell}$-Alg\`ebres, gradu\'ees respectivement par $\bN$ et $\bZ$.

\smallskip

\noindent
\underline{Exemple} : Le $\bZ_{\ell}$-faisceau $\bZ_{\ell}(1)$ d\'efini en (1.2.2) est inversible, et on conviendra dans la suite de poser :
\[
\bZ_{\ell}(r) = (\bZ_{\ell}(1))^{\otimes r} \qquad (r \in \bZ)\,.\hspace*{2cm}
\]

\bigskip

\noindent
\textbf{1.4.} \underline{Diverses cat\'egories construites \`a partir de la cat\'egorie des $\bZ_{\ell}$-faisceaux constructibles.}

\noindent
\textbf{1.4.1.} \underline{A-faisceaux constructibles} (A, $\bZ_{\ell}$-alg\`ebre finie).

Soit A une $\bZ_{\ell}$-alg\`ebre finie. On d\'efinit une cat\'egorie not\'ee
\[
\text{A-fc}(X)\,,\hspace*{5cm}
\]
dont les objets sont appel\'es A-\underline{faisceaux constructibles}, de la mani\`ere suivante.
\begin{itemize}
\item[(i)] Un objet de A-fc$(X)$ est un couple $(\cF, u)$ form\'e d'un $\bZ_{\ell}$-faisceau constructible $\cF$ et d'un morphisme de $\bZ_{\ell}$-alg\`ebres :
\[
u : A \longrightarrow \End(\cF)\,.\hspace*{2cm}
\]
\item[(ii)] Si $(\cF, u)$ et $(\cF', u')$ sont deux A-faisceaux constructibles, une fl\`eche $(\cF, u) \rightarrow (\cF', u')$ est un morphisme de $\bZ_{\ell}$-faisceaux
\[
f : \cF \longrightarrow \cF'\hspace*{3cm}
\]
tel que pour tout \'el\'ement $a$ de $A$ le diagramme ci-dessous soit commutatif :
\[
\begin{array}{ccc}
\cF & \xrightarrow{\;f\;} & \cF' \\
\downarrow{u(a)} & & \downarrow{u'(a)} \\
\cF & \xrightarrow{\;f\;} & \cF'
\end{array}
\]
\end{itemize}

\smallskip

\noindent
\textbf{1.4.2.} La cat\'egorie A-fc$(X)$ est une cat\'egorie ab\'elienne. Le foncteur oubli
\[
\text{A-fc}(X) \longrightarrow \bZ_{\ell}\text{-fc}(X)\hspace*{2cm}
\]
est exact et fid\`ele.

\smallskip

\noindent
\textbf{1.4.3.} Soit $\cF$ un A-faisceau constructible. On d\'efinit le \underline{dual} de $\cF$ par
\[
\cF^{\vee} = \cH om(\cF, A_X)\,.\hspace*{3cm}
\]

\smallskip

\noindent
\textbf{1.4.4.} Soit $f : X \rightarrow Y$ un morphisme de sch\'emas. Les foncteurs image r\'eciproque et image directe se prolongent de fa\c con \'evidente en des foncteurs entre les cat\'egories de A-faisceaux constructibles.

\smallskip

\noindent
\textbf{1.4.5.} Soit $\cF$ un $\bZ_{\ell}$-faisceau constructible, et notons $p$ le foncteur canonique $\bZ_{\ell}\text{-fc}(X) \rightarrow \bQ_{\ell}\text{-fc}(X)$. Le foncteur compos\'e de $p$ et du foncteur $\cG \mapsto \cF \otimes \cG$ rend inversibles les multiplications par des \'el\'ements de $\bZ_{\ell}$ et d\'efinit par suite un foncteur
\[
\bQ_{\ell}\text{-fc}(X) \longrightarrow \bQ_{\ell}\text{-fc}(X)\,,\hspace*{2cm}
\]
qui est exact \`a droite. Appliquant le m\^eme raisonnement au premier argument, on d\'efinit ainsi un bifoncteur exact \`a droite, appel\'e \underline{produit tensoriel} :
\[
\bQ_{\ell}\text{-fc}(X) \times \bQ_{\ell}\text{-fc}(X) \longrightarrow \bQ_{\ell}\text{-fc}(X)\,.\hspace*{1cm}
\]

Si $X$ est connexe et $\cF$ et $\cG$ sont deux $\bQ_{\ell}$-faisceaux constructibles constants tordus correspondant \`a des repr\'esentations de $\pi_1(X)$ sur les espaces vectoriels $V$ et $W$ respectivement, le produit tensoriel $\cF \otimes \cG$ est encore localement constant et correspond \`a la repr\'esentation diagonale de $\pi_1(X)$ sur $V \otimes W$.

\newpage

\noindent
\textbf{1.5.} \underline{Faisceaux $\ell$-adiques constructibles sur un spectre de corps.}

\noindent
\textbf{1.5.1.} Soit $k$ un corps s\'eparablement clos. Alors la cat\'egorie $\bZ_{\ell}\text{-fc}(\Spec(k))$ s'identifie \`a la cat\'egorie des $\bZ_{\ell}$-modules de type fini. La cat\'egorie $\bQ_{\ell}\text{-fc}(\Spec(k))$ s'identifie \`a celle des $\bQ_{\ell}$-espaces vectoriels de dimension finie.

\smallskip

\noindent
\textbf{1.5.2.} Soit $k$ un corps s\'eparablement clos et $G = \Gal(k_s/k)$ son groupe de Galois absolu. Alors la cat\'egorie $\bZ_{\ell}\text{-fc}(\Spec(k))$ s'identifie \`a la cat\'egorie des $\bZ_{\ell}$-modules de type fini munis d'une action continue de $G$. La cat\'egorie $\bQ_{\ell}\text{-fc}(\Spec(k))$ s'identifie \`a celle des repr\'esentations $\ell$-adiques continues de $G$, i.e. des $\bQ_{\ell}$-espaces vectoriels de dimension finie munis d'une action continue de $G$.

\smallskip

\noindent
\textbf{1.5.3.} Soient $k$ un corps et $\overline{k}$ une cl\^oture s\'eparable de $k$. Le foncteur ``fibre g\'eom\'etrique en le point g\'eom\'etrique $\Spec(\overline{k}) \rightarrow \Spec(k)$'' d\'efinit une \'equivalence de cat\'egories :
\[
\bZ_{\ell}\text{-fc}(\Spec(k)) \xrightarrow{\;\sim\;} \{\bZ_{\ell}\text{-modules de type fini munis d'une action continue de } \Gal(\overline{k}/k)\}\,.\hspace*{0.5cm}
\]

\smallskip

\noindent
\textbf{1.5.4.} Soit $k$ un corps fini \`a $q$ \'el\'ements. On d\'esigne par
\[
F_k \in \Gal(\overline{k}/k)\hspace*{4cm}
\]
le \'el\'ement de Frobenius g\'eom\'etrique. Une repr\'esentation $\ell$-adique de $\Gal(\overline{k}/k)$ est d\'etermin\'ee par la donn\'ee de l'action de $F_k$. Si $V$ est un $\bQ_{\ell}$-espace vectoriel de dimension finie muni d'une action continue de $\Gal(\overline{k}/k)$, les valeurs propres de $F_k$ agissant sur $V$ sont des nombres alg\'ebriques dont les inverses sont entiers sur $\bZ$.

\smallskip

\noindent
\textbf{1.5.5.} Soient $k$ un corps et $\overline{k}$ une cl\^oture s\'eparable. On d\'esigne par
\[
\bZ_{\ell}\text{-}\mathrm{modn}(\pi_1)\hspace*{4cm}
\]
la cat\'egorie des $\bZ_{\ell}$-modules de type fini munis d'une op\'eration continue de $\pi_1$.

\smallskip

\noindent
\textbf{LEMME 1.2.4.2.} Le foncteur pr\'ec\'edent
\[
\varprojlim_{\bN} : (\ell\text{-}\bZ)\text{-ad}(\mathrm{Abf}(\pi_1)) \longrightarrow \bZ_{\ell}\text{-}\mathrm{modn}(\pi_1)
\]
est une \underline{\'equivalence de cat\'egories.}

\smallskip

\noindent
\underline{Preuve} : Il est clair qu'il est g\'en\'eriquement surjectif : un $\bZ_{\ell}$-module de type fini $E$ sur lequel $\pi_1$ op\`ere contin\^ument est l'image du syst\`eme projectif
\[
(E/\ell^{n+1}E)_{n \in \bN}\,,\hspace*{4cm}
\]
avec les op\'erations \'evidentes de $\pi_1$.

\newpage

\noindent
\textbf{2.} \underline{Formalism de la cohomologie $\ell$-adique.}

\noindent
\textbf{2.1.} \underline{Image r\'eciproque.}

Soit $f : X \rightarrow Y$ un morphisme de pr\'esch\'emas localement noeth\'eriens. Le foncteur image r\'eciproque
\[
f^* : \mathrm{Ab}(Y) \longrightarrow \mathrm{Ab}(X)\hspace*{2cm}
\]
est exact et se prolonge donc de fa\c con \'evidente en un foncteur \underline{exact}, not\'e de la m\^eme mani\`ere,
\[
f^* : \underline{\Hom}(\bN^{\circ}, \mathrm{Ab}(Y)) \longrightarrow \underline{\Hom}(\bN^{\circ}, \mathrm{Ab}(X))\,.\hspace*{1cm}
\]

Ce dernier foncteur transforme \'evidemment $\bZ_{\ell}$-faisceau constructible en $\bZ_{\ell}$-faisceau constructible, et syst\`eme projectif localement AR-nul en syst\`eme projectif localement AR-nul. Avec la notation de (1.5.4), il induit donc deux foncteurs ``images r\'eciproques''
\[
\begin{array}{rcl}
f^* : \bZ_{\ell}\text{-fc}(Y) & \longrightarrow & \bZ_{\ell}\text{-fc}(X)\,, \\
f^* : \mathrm{U}(Y) & \longrightarrow & \mathrm{U}(X)\,,\hspace*{1cm}
\end{array}
\]
le second \'etant exact.

L'\'equivalence de cat\'egories (1.5.5) et le fait que $f^*$ transforme syst\`eme projectif localement AR-nul en syst\'eme projectif localement AR-nul montre alors que le foncteur image r\'eciproque
\[
f^* : \bZ_{\ell}\text{-fc}(Y) \longrightarrow \bZ_{\ell}\text{-fc}(X)\hspace*{1cm}
\]
est aussi exact.

Si $g : Y \rightarrow Z$ est un autre morphisme de pr\'esch\'emas localement noeth\'eriens, les foncteurs $f^*$ et $g^*$ sont \'evidemment li\'es par un isomorphisme
\[
(g \circ f)^* \xrightarrow{\;\sim\;} f^* \circ g^*\,,\hspace*{3cm}
\]
avec la propri\'et\'e d'``associativit\'e'' habituelle (SGA 1 VI 7 et suiv.)

\bigskip

\noindent
\textbf{2.2.} \underline{Suite spectrale, images directes sup\'erieures.}

\noindent
\textbf{2.2.1.} Soient $f : X \rightarrow Y$ un morphisme compactifiable de pr\'esch\'emas localement noeth\'eriens, et $\cF$ un $\bZ_{\ell}$-faisceau constructible sur $X$. On d\'efinit les images directes sup\'erieures
\[
R^i f_!(\cF)\hspace*{5cm}
\]
comme les objets de $\bZ_{\ell}\text{-fc}(Y)$ d\'efinis par passage \`a la limite :
\[
R^i f_!(\cF) = (R^i f_!(\cF_n))_{n \in \bN}\,.\hspace*{3cm}
\]

\smallskip

\noindent
\textbf{2.2.2.} \underline{Suite spectrale de Leray.} Soient $f : X \rightarrow Y$ et $g : Y \rightarrow Z$ deux morphismes compactifiables de pr\'esch\'emas localement noeth\'eriens, et $\cF$ un $\bZ_{\ell}$-faisceau constructible sur $X$. Il existe une suite spectrale
\[
E_2^{pq} = R^p g_! (R^q f_! (\cF)) \Rightarrow R^{p+q}(g \circ f)_! (\cF)\,,\hspace*{1cm}
\]
functorielle en $\cF$.

\smallskip

\noindent
\textbf{2.2.3.} \underline{Changement de base.} Consid\'erons un diagramme cart\'esien
\[
\begin{array}{ccc}
X' & \xrightarrow{\;g'\;} & X \\
\downarrow{f'} & & \downarrow{f} \\
Y' & \xrightarrow{\;g\;} & Y
\end{array}
\]
o\`u $f$ (donc $f'$) est compactifiable et les sch\'emas sont localement noeth\'eriens. Pour tout $\bZ_{\ell}$-faisceau constructible $\cF$ sur $X$, il existe des morphismes de changement de base
\[
g^* (R^i f_! (\cF)) \longrightarrow R^i f'_! (g'^* \cF)\,,\hspace*{2cm}
\]
qui sont des isomorphismes lorsque $g$ est propre ou $f$ est lisse.

\smallskip

\noindent
\textbf{2.2.4.} \underline{Suite exacte longue de cohomologie.} Soient $f : X \rightarrow Y$ un morphisme compactifiable et
\[
0 \longrightarrow \cF' \longrightarrow \cF \longrightarrow \cF'' \longrightarrow 0\hspace*{1cm}
\]
une suite exacte de $\bZ_{\ell}$-faisceaux constructibles sur $X$. Alors il existe une suite exacte longue
\[
\cdots \longrightarrow R^i f_!(\cF') \longrightarrow R^i f_!(\cF) \longrightarrow R^i f_!(\cF'') \longrightarrow R^{i+1} f_!(\cF') \longrightarrow \cdots\hspace*{0.5cm}
\]

\smallskip

\noindent
\textbf{2.2.5.} Le formalisme d\'evelopp\'e dans ce paragraphe pour les $\bZ_{\ell}$-faisceaux s'\'etend sans difficult\'e aux A-faisceaux ($A$ une $\bZ_{\ell}$-alg\`ebre finie) et aux L-faisceaux ($L$ extension finie de $\bQ_{\ell}$). En ce qui concerne par exemple le cas des $\bQ_{\ell}$-faisceaux, le point fondamental pour les d\'emonstrations est que

\newpage

les foncteurs $R^i f_!$ commutent au passage au quotient par la sous-cat\'egorie des faisceaux de torsion. Cela r\'esulte du fait que, pour tout entier $n$, le foncteur $R^i f_!$ sur les faisceaux de $\bZ/\ell^{n+1}\bZ$-modules est exact \`a droite et commute aux limites inductives filtrantes.

\bigskip

\noindent
\textbf{2.3.} \underline{Cohomologie des fibr\'es projectifs.}

\noindent
\textbf{2.3.1.} Soient $S$ un sch\'ema, $E$ un $\cO_S$-Module localement libre de rang $r+1$, et
\[
p : \mathbb{P}(E) \longrightarrow S\hspace*{4cm}
\]
le fibr\'e projectif associ\'e. On pose $X = \mathbb{P}(E)$. On note
\[
\xi = c_1(\cO_X(1)) \in H^2(X, \bZ_{\ell}(1))\hspace*{3cm}
\]
la premi\`ere classe de Chern du faisceau inversible standard $\cO_X(1)$. Pour tout entier $i \geq 0$, on a une classe de cohomologie $\xi^i \in H^{2i}(X, \bZ_{\ell}(i))$.

\smallskip

\noindent
\textbf{THEOREME 2.3.2.} \underline{Les morphismes}
\[
\theta_i : H^{k-2i}(S, \bZ_{\ell}(-i)) \longrightarrow H^k(X, \bZ_{\ell})\,, \quad x \longmapsto p^*(x) \cdot \xi^i\,,
\]
\underline{d\'efinissent un isomorphisme}
\[
\bigoplus_{i=0}^{r} H^{k-2i}(S, \bZ_{\ell}(-i)) \xrightarrow{\;\sim\;} H^k(X, \bZ_{\ell})\,.\hspace*{1cm}
\]

\smallskip

\noindent
\underline{Preuve} : On proc\'ede par r\'ecurrence sur $r$. Pour $r = 0$, l'assertion est triviale. Supposons le th\'eor\`eme d\'emontr\'e pour les fibr\'es projectifs de rang $\leq r-1$. On utilise la suite exacte de localisation pour l'ouvert compl\'ementaire d'une section hyperplane. Soient $E'$ un sous-module localement facteur direct de rang $r$ de $E$, et $H = \mathbb{P}(E')$ l'hyperplane correspondant. La suite exacte de cohomologie \`a supports propres pour l'ouvert $U = X \setminus H$ donne :
\[
\cdots \longrightarrow H^{k-2}(H, \bZ_{\ell}(-1)) \longrightarrow H^k(X, \bZ_{\ell}) \longrightarrow H^k(U, \bZ_{\ell}) \longrightarrow H^{k-1}(H, \bZ_{\ell}(-1)) \longrightarrow \cdots\hspace*{0.5cm}
\]

Par l'hypoth\`ese de r\'ecurrence, on conna\^it la cohomologie de $H$. D'autre part, $U$ est un espace affine sur $S$, donc $H^k(U, \bZ_{\ell}) = H^k(S, \bZ_{\ell})$ pour tout $k$. Un calcul direct permet alors de conclure.

\smallskip

\noindent
\textbf{2.3.3.} \underline{Structure de l'anneau de cohomologie.} Le th\'eor\`eme 2.3.2 montre que le $H^*(S, \bZ_{\ell})$-module $H^*(X, \bZ_{\ell})$ est libre de base $1, \xi, \xi^2, \ldots, \xi^r$. L'anneau $H^*(X, \bZ_{\ell})$ est donc isomorphe \`a
\[
H^*(S, \bZ_{\ell})[T]/(T^{r+1})\,,\hspace*{3cm}
\]
o\`u $T$ correspond \`a $\xi \in H^2(X, \bZ_{\ell}(1))$.

\newpage

\noindent
\textbf{2.4.} \underline{Th\'eor\`eme de changement de base propre.}

\noindent
\textbf{THEOREME 2.4.1.} Soient
\[
\begin{array}{ccc}
X' & \xrightarrow{\;g'\;} & X \\
\downarrow{f'} & & \downarrow{f} \\
Y' & \xrightarrow{\;g\;} & Y
\end{array}
\]
un diagramme cart\'esien de sch\'emas localement noeth\'eriens, avec $f$ propre. Soit $\cF$ un $\bZ_{\ell}$-faisceau constructible sur $X$. Alors les morphismes de changement de base
\[
g^*(R^i f_* \cF) \longrightarrow R^i f'_*(g'^* \cF)\hspace*{2cm}
\]
\underline{sont des isomorphismes.}

\smallskip

\noindent
\underline{Preuve} : On se ram\`ene aussit\^ot au cas o\`u $Y$ et $Y'$ sont strictement locaux, $g$ \'etant le morphisme canonique du point g\'eom\'etrique ferm\'e. L'assertion r\'esulte alors du th\'eor\`eme de changement de base propre pour la cohomologie \`a supports propres (SGA 4 XII 5.1) et du fait que, $f$ \'etant propre, on a $R^i f_! = R^i f_*$.

\bigskip

\noindent
\textbf{2.5.} \underline{Th\'eor\`eme de finitude.}

\noindent
\textbf{THEOREME 2.5.1.} Soient $S$ un sch\'ema r\'egulier de type fini sur un corps $k$ et $f : X \rightarrow S$ un morphisme propre. Soit $\cF$ un $\bZ_{\ell}$-faisceau constructible sur $X$. Alors les $R^i f_* \cF$ sont des $\bZ_{\ell}$-faisceaux constructibles sur $S$.

\smallskip

\noindent
\underline{Preuve} : Compte tenu du th\'eor\`eme de changement de base propre 2.4.1, il suffit de montrer que pour tout point g\'eom\'etrique $\overline{s}$ de $S$, les $\bZ_{\ell}$-modules $H^i(X_{\overline{s}}, \cF)$ sont de type fini. On se ram\`ene ainsi au cas o\`u $S$ est le spectre d'un corps s\'eparablement clos. L'assertion r\'esulte alors du th\'eor\`eme de finitude pour la cohomologie des sch\'emas propres sur un corps (SGA 4 XIV 1.1).

\bigskip

\noindent
\textbf{2.6.} \underline{Cohomologie \`a supports propres.}

\noindent
\textbf{2.6.1.} Soient $f : X \rightarrow S$ un morphisme s\'epar\'e de type fini de sch\'emas localement noeth\'eriens, et $\cF$ un $\bZ_{\ell}$-faisceau constructible sur $X$. On d\'efinit la cohomologie \`a supports propres
\[
H^i_c(X/S, \cF)\hspace*{4cm}
\]
comme le $\bZ_{\ell}$-faisceau constructible sur $S$ d\'efini par compactification de $f$.

\smallskip

\noindent
\textbf{2.6.2.} Soient $j : U \rightarrow X$ une immersion ouverte et $i : Z \rightarrow X$ l'immersion ferm\'ee compl\'ementaire. Pour tout $\bZ_{\ell}$-faisceau constructible $\cF$ sur $U$, on a une suite exacte (fonctorielle en $\cF$)
\[
\cdots \longrightarrow R^i f_!(j_! \cF) \longrightarrow R^i f_!(Rj_* \cF) \longrightarrow R^i f_!(i_* i^* Rj_* \cF) \longrightarrow \cdots\hspace*{0.5cm}
\]

\bigskip

\noindent
\textbf{3.} \underline{Fonctions} $L$ \underline{et} $Z$ \underline{$\ell$-adiques.}

\noindent
\textbf{3.1.} \underline{D\'efinitions.}

\noindent
\textbf{3.1.1.} Soient $\mathbb{F}_q$ un corps fini \`a $q$ \'el\'ements et $X$ un sch\'ema de type fini sur $\mathbb{F}_q$. On d\'esigne par $|X|$ l'ensemble des points ferm\'es de $X$. Pour tout $x \in |X|$, on note $d(x)$ le degr\'e du corps r\'esiduel $k(x)$ sur $\mathbb{F}_q$, et $N(x) = q^{d(x)}$ le nombre d'\'el\'ements de $k(x)$.

\smallskip

\noindent
\textbf{3.1.2.} Soit $\cF$ un $\bQ_{\ell}$-faisceau constructible sur $X$. Pour tout $x \in |X|$, la fibre g\'eom\'etrique $\cF_{\overline{x}}$ est un $\bQ_{\ell}$-espace vectoriel de dimension finie, muni d'une action continue du groupe de Galois $\Gal(\overline{k(x)}/k(x))$. L'\'el\'ement de Frobenius g\'eom\'etrique $F_x$ agit sur $\cF_{\overline{x}}$ avec des valeurs propres $\alpha_1, \ldots, \alpha_n$ qui sont des nombres alg\'ebriques. On pose
\[
\det(1 - F_x t^{d(x)} \mid \cF_{\overline{x}}) = \prod_{i=1}^{n}(1 - \alpha_i t^{d(x)})\,.\hspace*{1cm}
\]

\smallskip

\noindent
\textbf{3.1.3.} On d\'efinit la fonction $Z$ associ\'ee \`a $\cF$ par :
\[
Z(X, \cF; t) = \prod_{x \in |X|} \frac{1}{\det(1 - F_x t^{d(x)} \mid \cF_{\overline{x}})} \in \bQ_{\ell}[[t]]\,.\hspace*{0.5cm}
\]
Lorsque $\cF = \bQ_{\ell}$, on note simplement $Z(X; t)$ la fonction $Z$ correspondante.

\smallskip

\noindent
\textbf{3.2.} \underline{Formule de Lefschetz.}

\noindent
\textbf{THEOREME 3.2.1.} \underline{(Formule des traces de Lefschetz).} Soient $X$ un sch\'ema propre et lisse sur $\mathbb{F}_q$ et $\cF$ un $\bQ_{\ell}$-faisceau constructible. Alors :
\[
Z(X, \cF; t) = \prod_{i=0}^{2\dim X} \det(1 - F^* t \mid H^i(X \otimes \overline{\mathbb{F}}_q, \cF))^{(-1)^{i+1}}\,.\hspace*{0.5cm}
\]

\smallskip

\noindent
\underline{Preuve} : La preuve utilise la formule de Lefschetz pour les correspondances sur les vari\'et\'es propres et lisses (SGA 5 III), ainsi que la th\'eorie du \textsl{foncteur $L$} d\'Artin-Verdier.

\bigskip

\noindent
\textbf{3.3.} \underline{\'Enonc\'e des conjectures de Weil.}

\noindent
\textbf{3.3.1.} Soient $X$ un sch\'ema propre et lisse sur un corps fini $\mathbb{F}_q$, et $\ell$ un nombre premier diff\'erent de la caract\'eristique de $\mathbb{F}_q$. Les conjectures de Weil affirment que :
\begin{itemize}
\item[(i)] Les polyn\^omes caract\'eristiques $\det(1 - F^* t \mid H^i(X \otimes \overline{\mathbb{F}}_q, \bQ_{\ell}))$ sont \`a coefficients entiers, ind\'ependants de $\ell$.
\item[(ii)] Les valeurs propres de $F^*$ agissant sur $H^i(X \otimes \overline{\mathbb{F}}_q, \bQ_{\ell})$ sont des entiers alg\'ebriques de valeur absolue complexe $q^{i/2}$.
\item[(iii)] (Dualit\'e de Poincar\'e) Il existe une dualit\'e parfaite
\[
H^i(X \otimes \overline{\mathbb{F}}_q, \bQ_{\ell}) \times H^{2d-i}(X \otimes \overline{\mathbb{F}}_q, \bQ_{\ell}(d)) \longrightarrow \bQ_{\ell}\,.\hspace*{0.5cm}
\]
\end{itemize}

\smallskip

\noindent
\textbf{3.3.2.} Ces conjectures ont \'et\'e prouv\'ees par Deligne (SGA 4 $\frac{1}{2}$, SGA 5). La preuve utilise de fa\c con essentielle la formule des traces et la th\'eorie des poids de Hodge-Deligne.

\newpage

\noindent
\textbf{4.} \underline{Cohomologie des compl\'et\'es $\ell$-adiques.}

\noindent
\textbf{4.1.} \underline{Passage \`a la limite projective.}

\noindent
\textbf{4.1.1.} Soit $X$ un sch\'ema propre sur $\Spec(\bZ[1/\ell])$. Pour tout entier $n \geq 1$, on pose $X_n = X \otimes_{\bZ[1/\ell]} \bZ[1/\ell, \zeta_{\ell^n}]$, o\`u $\zeta_{\ell^n}$ est une racine primitive $\ell^n$-\`eme de l'unit\'e. Les groupes de cohomologie
\[
H^i(X_n, \bZ/\ell^n\bZ)\hspace*{3cm}
\]
forment un syst\`eme projectif dont la limite projective
\[
H^i_{\ell\text{-}\mathrm{comp}}(X) = \varprojlim_{n} H^i(X_n, \bZ/\ell^n\bZ)\hspace*{1cm}
\]
est appel\'ee la $i$-\`eme groupe de cohomologie compl\'et\'ee $\ell$-adique de $X$.

\smallskip

\noindent
\textbf{4.1.2.} Les groupes $H^i_{\ell\text{-}\mathrm{comp}}(X)$ sont des $\bZ_{\ell}$-modules de type fini. Ils co\"  ncident avec les groupes de cohomologie $\ell$-adique d\'efinis pr\'ec\'edemment pour les sch\'emas propres et lisses.

\smallskip

\noindent
\textbf{4.1.3.} Soient $X$ un sch\'ema propre sur $\Spec(\bZ[1/\ell])$ et $\cF$ un $\bZ_{\ell}$-faisceau constructible sur $X$. On d\'efinit de mani\`ere analogue les groupes de cohomologie compl\'et\'ee
\[
H^i_{\ell\text{-}\mathrm{comp}}(X, \cF) = \varprojlim_{n} H^i(X_n, \cF/\ell^n \cF)\,.\hspace*{1cm}
\]

\bigskip

\noindent
\textbf{4.2.} \underline{Suite spectrale d'Iwasawa.}

\noindent
\textbf{4.2.1.} Soient $X$ un sch\'ema propre sur $\bZ[1/\ell]$ et $G = \Gal(\bQ(\zeta_{\ell^{\infty}})/\bQ) \cong \bZ_{\ell}^{\times}$. Le groupe $G$ op\`ere sur les groupes $H^i_{\ell\text{-}\mathrm{comp}}(X)$ par transport de structure. On peut donc consid\'erer ces groupes comme des modules sur l'alg\'ebre d'Iwasawa
\[
\Lambda = \bZ_{\ell}[[G]] = \varprojlim_{n} \bZ_{\ell}[G/\ell^n \bZ_{\ell}^{\times}]\,.\hspace*{1cm}
\]

\smallskip

\noindent
\textbf{4.2.2.} \underline{Th\'eor\`eme.} Les $\Lambda$-modules $H^i_{\ell\text{-}\mathrm{comp}}(X)$ sont de type fini.

\smallskip

\noindent
\underline{Preuve} : La preuve utilise la suite spectrale de Hochschild-Serre pour le rev\^etement $X_n \rightarrow X$ et le th\'eor\`eme de finitude pour la cohomologie des sch\'emas propres.

\bigskip

\noindent
\textbf{5.} \underline{Applications.}

\noindent
\textbf{5.1.} \underline{Th\'eor\`eme de rationalit\'e des fonctions $Z$.}

\noindent
\textbf{THEOREME 5.1.1.} Soient $X$ un sch\'ema de type fini sur $\mathbb{F}_q$ et $\cF$ un $\bQ_{\ell}$-faisceau constructible sur $X$. Alors la fonction $Z(X, \cF; t)$ est une fraction rationnelle en $t$.

\smallskip

\noindent
\underline{Preuve} : Lorsque $X$ est propre et lisse, cela r\'esulte de la formule des traces de Lefschetz 3.2.1. Le cas g\'en\'eral s'en d\'eduit par d\'evissage \`a la d\'ecomposition cellulaire et la suite exacte longue de cohomologie.

\bigskip

\noindent
\textbf{5.2.} \underline{Th\'eor\`eme de l'unit\'e.}

\noindent
\textbf{THEOREME 5.2.1.} Soient $K$ un corps de nombres, $\cO_K$ l'anneau de ses entiers, et $\ell$ un nombre premier. Alors le groupe des unit\'es $\cO_K^{\times}$ est de type fini, et son rang est donn\'e par
\[
\rg(\cO_K^{\times}) = r_1 + r_2 - 1\,,\hspace*{4cm}
\]
o\`u $r_1$ est le nombre de places r\'eelles et $r_2$ le nombre de places complexes de $K$.

\smallskip

\noindent
\underline{Preuve} : Ce th\'eor\`eme classique peut \^etre interpr\'et\'e en termes de cohomologie $\ell$-adique. La suite spectrale d'Iwasawa pour $\Spec(\cO_K[1/\ell])$ fournit une description du groupe des unit\'es en termes de modules d'Iwasawa.

\bigskip

\noindent
\textbf{5.3.} \underline{Structure galoisienne de la cohomologie.}

\noindent
\textbf{THEOREME 5.3.1.} Soient $k$ un corps de nombres et $X$ une vari\'et\'e projective et lisse sur $k$. Soit $\overline{k}$ une cl\^oture alg\'ebrique de $k$. Alors la repr\'esentation de $\Gal(\overline{k}/k)$ sur $H^i_{\text{\'et}}(X \otimes_k \overline{k}, \bQ_{\ell})$ est semi-simple, et ses facteurs de composition sont ind\'ependants de $\ell$ (pour $\ell$ diff\'erent de la caract\'eristique r\'esiduelle).

\smallskip

\noindent
\underline{Preuve} : La semi-simplicit\'e r\'esulte de la th\'eorie des poids de Deligne. L'ind\'ependance en $\ell$ est une cons\'equence de la conjecture de Weil prouv\'ee par Deligne et de la th\'eorie des cycles alg\'ebriques.

\newpage

\begin{center}
\textit{FIN DE L'EXPOSE VI}
\end{center}
